% 本文件由 LaTeX 排版 Agent 根据有效 translation.json 自动生成。
% author=Athanasius; work_id=2815; title=亚略派史

\chapter{亚略派史}

% structure: p0001 source=h1 target=omitted reason=page-title-already-rendered-by-parent

第一部 \Person{君士坦丁}时期的亚略派迫害。\newline{}第二部 \Person{君士坦提乌斯}时期的首次亚略派迫害。\newline{}第三部 撒狄卡会议上天主教徒的恢复。\newline{}第四部 \Person{君士坦提乌斯}时期的第二次亚略派迫害。\newline{}第五部 \Person{利伯略}的迫害与背教。\newline{}第六部 \Person{霍西乌斯}的迫害与背教。\newline{}第七部 \Place{亚历山大里亚}的迫害。\newline{}第八部 \Place{埃及}的迫害。\par

\SourceRecord{Translated by M. Atkinson and Archibald Robertson.；Nicene and Post-Nicene Fathers, Second Series；Vol. 4.；Edited by Philip Schaff and Henry Wace.；Buffalo, NY: Christian Literature Publishing Co.,；1892.；<http://www.newadvent.org/fathers/2815.htm>.}

% structure: p0001 source=h1 target=section reason=page-title

\section{亚略派历史，第一部分}

\Signature{圣亚大纳削}\par

% structure: p0003 source=h2 target=subsection reason=relative-heading-rank; base=h2

\subsection{一、\Person{君士坦丁}治下的亚略派迫害}

不久之后，他们便实施了其计谋，正是为了这些计谋，他们才采用了这些诡计；因为他们一形成计划，就立刻接纳了\Person{亚略}及其同党共融。他们搁置了对他们一再作出的谴责，并再次假借皇帝的权威为他们撑腰。他们在信中竟不知羞耻地说：\Quote{既然\Person{亚大纳削}遭受了苦难，一切嫉妒都已消除，今后我们就接纳\Person{亚略}及其同党吧}；为了恐吓听者，又补充说：\Quote{因为皇帝已下了命令。}此外，他们还不以为耻地补充说：\Quote{因为这些人宣认正统信仰}；毫不畏惧经上所载的：\BibleQuote{祸哉，那些称苦为甜，以暗为光的人！}\BibleRef{依 5:20}，因为他们随时准备不择手段来支持他们的异端。现在，难道这还不是向所有人清楚地证明，我们从前所受的苦，以及你们如今对我们的迫害，并非依据教会判决的权威，而是基于皇帝的威胁，以及对基督的虔诚吗？同样，他们也以类似的方式阴谋对付其他主教，同样捏造指控来陷害他们；其中一些人在流放之地安眠，获得了基督信仰宣认的荣耀；还有一些人仍被驱逐出自己的祖国，并越来越英勇地抵抗他们的异端，说：\BibleQuote{谁能使我们与基督的爱隔绝呢？}\BibleRef{罗 8:35}。\par

% structure: p0005 source=h2 target=subsection reason=relative-heading-rank; base=h2

\subsection{二、亚略派为党派牺牲道德与正直}

由此你也可以辨别它的性质，并能更坚决地谴责它。一个人，只要他是他们的朋友，与他们一同不虔敬，即使他因所犯的无数其他恶行而受到万千指控；即使对他不利的证据和证明再清楚不过；他们也会赞许他，他立刻便成为皇帝的朋友，凭借他的不虔敬而获得引荐；并且制造许多借口，从而在官府面前胆敢为所欲为。但是，那揭露他们的不虔敬，并诚实地为基督的事业辩护的人，即使他在各方面纯洁无瑕，即使他自觉毫无过失，即使无人控告他；然而，基于他们为他捏造的虚假借口，他会立即被抓捕，并由皇帝下令处以流放，仿佛他确实犯有他们想要加给他的罪行，或者像\Person{纳波特}那样，仿佛他侮辱了国王；而那为他们的异端事业辩护的人，则被寻找，并立即被派去占据另一人的教会；从此以后，凡不接纳他的人，就要遭受没收财产、侮辱和各种残酷的对待。而最为奇怪的是，那人民所渴望、并知道是无可指摘的人\BibleRef{弟前 3:2}，皇帝却带走并流放；而那人民既不渴望、也不认识的人，他却从远方派来，带着士兵和他自己的信函。从此，他们便被迫做出艰难的抉择：要么恨他们原来所爱的那位，他曾是他们的导师和虔敬上的父亲；而去爱他们所不渴望的人，并将自己的子女托付给一个他们对其生活、品行和人格一无所知的人；要么，如果他们不服从皇帝，就必定要遭受惩罚。\par

% structure: p0007 source=h2 target=subsection reason=relative-heading-rank; base=h2

\subsection{三、他们行事的肆无忌惮}

不虔敬者现今就是这样一如既往地攻击正统信徒，在各地众人面前显明他们的恶意与不虔敬。姑且承认他们指控了\Person{亚大纳削}；但其他主教又做了什么？他们凭什么指控他们呢？在他们那里也找到了\Person{阿尔塞尼乌斯}的尸体吗？有长老\Person{马卡里乌斯}，或者他们中间有\OriginalTerm{圣爵}{cup}被打破吗？有\OriginalTerm{梅勒提乌斯派}{Meletian}的人假作伪善吗？不，但是，正如他们对其他主教的所作所为表明，他们对\Person{亚大纳削}的指控极有可能是虚假的；同样，他们对\Person{亚大纳削}的攻击也清楚表明，他们对其他主教的指控同样毫无根据。这一异端如同某个巨大的怪物来到世间，不仅用其言语如同用牙齿伤害无辜者，而且还雇用了外在的权力来辅助其阴谋。奇怪的是，正如我先前所说，他们中无人受到指控；或者即使有人被指控，他也不被审判；或者即便只是做出调查的样子，他也会在证据面前被宣告无罪，而定罪的一方反而遭到阴谋陷害，而不是罪犯蒙羞。因此，他们整个团伙充满了闲散；而他们的密探——因为他们根本不是主教——是他们中最卑鄙的。倘若他们之中有谁想成为主教，听到的不是\BibleQuote{主教必须无可指摘}\BibleRef{弟前 3:2}，而只是：\Quote{采取与基督相反的意见，不必拘泥于品行。这就足以使你获得宠幸，并与皇帝建立友谊。}这就是那些支持\Person{亚略}主张之人的品行。而那些热心追求真理的人，无论他们如何彰显自己的圣善与纯洁，却依然如我前所说，只要这些人愿意，随便他们想捏造什么借口，随时都会变成罪犯。这一点，如我之前所提及的，你从他们的行事中便可清楚看出。\par

% structure: p0009 source=h2 target=subsection reason=relative-heading-rank; base=h2

\subsection{四、亚略派迫害\Person{欧斯塔修斯}等人}

有一位\Place{安提约基雅}主教\Person{欧斯塔提乌斯}，是一位\OriginalTerm{精修圣人}{Confessor}，\OriginalTerm{信德}{Faith}纯正。此人因为极其热心于真理，憎恶\OriginalTerm{亚略异端}{Arian heresy}，且不接纳那些采纳其信条的人，便在\Person{君士坦丁}皇帝面前被诬告，并被捏造了一个指控，说他侮辱了皇帝的母亲。他立刻被处以流放，与他一同被流放的还有大批长老和执事。主教被流放后，那些因不虔敬而不被他接纳进入\OriginalTerm{神职界}{clerical order}的人，不仅被他们接纳进教会，而且其中大部分人还被任命为主教，好让他们在不虔敬的事上有同伙。这些人中有现在在\Place{安提约基雅}的阉人\Person{莱昂提乌斯}，以及他的前任\Person{斯特法努斯}、\Place{劳狄刻雅}的\Person{乔治}、\Place{特里波利}的\Person{狄奥多西}、\Place{耳曼尼基亚}的\Person{欧多克西乌斯}，还有现在在\Place{瑟巴斯提亚}的\Person{欧斯塔提乌斯}。\par

5. 难道他们就到此为止了吗？不。因为\Place{亚德里亚堡}的主教\Person{欧特罗皮乌斯}，一个善良且在各方面都很优秀的人，只因常揭露\Person{欧瑟伯}，并劝诫那些路过的人不要遵从他不虔敬的指示，就遭受了与\Person{欧斯塔提乌斯}相同的待遇，被驱逐出自己的城市和教会。\Person{巴西利纳}在针对他的诉讼中最为积极。\Place{巴拉内亚}的\Person{欧弗拉提翁}、\Place{帕尔图斯}的\Person{基马提乌斯}、\Place{安塔拉杜斯}的\Person{卡尔特里乌斯}、\Place{加沙}的\Person{阿斯克勒帕斯}、\Place{叙利亚}的\Place{贝罗亚}的\Person{居鲁士}、\Place{亚细亚}的\Person{狄奥多鲁斯}、\Place{西尔米翁}的\Person{多姆尼翁}，以及\Place{特里波利}的\Person{埃拉尼库斯}，仅仅因被人知道憎恶\OriginalTerm{异端}{heresy}，就有的以这样那样的借口，有的毫无借口，便被他们凭皇帝书信的权威撤职，逐出他们的城市，并另立他们所知道的不虔敬之人占据他们的教会。\par

% structure: p0012 source=h2 target=subsection reason=relative-heading-rank; base=h2

\subsection{六、马尔切路斯案件}

关于\Place{加拉太}的主教\Person{马尔切路斯}，我或许无需多言；因为所有人都已听说，曾被\Person{马尔切路斯}率先指控不虔敬的\Person{欧瑟伯}及其同伙，如何反控他，并导致这老人被流放。他前往\Place{罗马}，在那里为自己辩护，并应他们的要求，提交了一份书面信德宣言，\OriginalTerm{撒狄卡会议}{Council of Sardica}对此表示认可。但\Person{欧瑟伯}及其同伙不但没有辩护，即便从他们的著作中被证实不虔敬，他们也不以为耻，反而对所有人愈加胆大妄为。因为他们经由那些妇女得以接近皇帝，便在众人面前令人生畏。\par

% structure: p0014 source=h2 target=subsection reason=relative-heading-rank; base=h2

\subsection{七、\Place{君士坦丁堡}的\Person{保禄}殉道}

我想，没有人不知道\Place{君士坦丁堡}的主教\Person{保禄}的案件；因为一座城市越是显赫，其中发生的事就越是无法隐瞒。针对他也捏造了一项指控。因为控告他的\Person{马切多尼乌斯}，就是现今接替他成为主教的那个人（我自己当时也在控告现场），后来却与他共融，并且曾是\Person{保禄}手下的一位长老。然而，当\Person{欧瑟伯}心怀嫉妒，企图夺取该城的主教职位时（他此前以同样的方式从\Place{贝里图斯}转任至\Place{尼科美底亚}），对\Person{保禄}的指控被重新提起；他们并未放弃阴谋，反而坚持诽谤。他先是被\Person{君士坦丁}流放到\Place{本都}；第二次，被\Person{君士坦提乌斯}下令戴上铁链，送往\Place{美索不达米亚}的\Place{辛加拉}；又从那里被转到\Place{埃美萨}；第四次，被流放到\Place{卡帕多细亚}的\Place{库库苏斯}，靠近\Place{塔鲁斯山}的荒野；在那里，据那些与他同在的人说，他被他们亲手勒死。然而，这些从不说真话的人，虽然犯下此罪，却在他死后还不羞耻地编造另一个说法，声称他是因病而死；尽管所有住在那地方的人都清楚实情。就连当时担任那一地区代理总督的\Person{菲拉格里乌斯}，虽然完全按照他们的意愿呈报所有事情，也对此事感到震惊；或许因干这恶事的是别人而非自己而苦恼，他便告诉主教\Person{塞拉皮翁}和我们的许多朋友，说是他们把\Person{保禄}关在一个非常狭窄黑暗的地方，任其饿死；六天后，他们进去发现他还活着，便立刻攻击此人，将他勒死。这就是他生命的终结；他们说，总督\Person{斐理伯}是这起谋杀的执行者。然而，\OriginalTerm{天主的公义}{Divine Justice}并没有忽略此事；因为不到一年，\Person{斐理伯}就被极大的耻辱剥夺了官职，以致沦为一介平民，成为那些他最不希望见到自己垮台之人的笑柄。他内心极度痛苦，像\Person{加音}一样呻吟战栗，天天等着有人来杀他；远离故土和亲友，他就在自己最不愿承受的境况中死去，仿佛被自己的厄运惊呆了。不但如此，这些人对他们在生前捏造指控的对象，即使死后也不放过。他们如此急于向所有的人显示自己的威风，以致流放活人，又不怜悯死者；反而在全世界人中，唯独他们对已亡者表现仇恨，并图谋反对其友人；这些真正不人道、憎恶良善、性情野蛮甚于仇敌的人，因为他们不虔敬的缘故，不顾真理，仅凭虚假的指控，就急切地谋划毁灭我和所有其余的人。\par

% structure: p0016 source=h2 target=subsection reason=relative-heading-rank; base=h2

\subsection{八、公教徒的恢复}

三位兄弟，\Person{君士坦丁}、\Person{君士坦提乌斯}和\Person{君士坦斯}，看到这种情况，便在父亲去世后，使所有的人返回自己的国家和教会；他们给各自的教会写了关于其他人的信，关于\Person{亚大纳削}，他们则写了如下文字；这同样表明了整个过程的暴力，并证实了\Person{欧瑟伯}及其同伙的杀人心性。\par

\Emph{\Person{君士坦丁凯撒}致\Place{亚历山大}城\OriginalTerm{天主教会}{Catholic Church}民众的一封信的副本。}\par

\Quote{我猜想，你们虔诚的心灵尚未忽略……等等。}\par

这就是他的信；对于他们的阴谋，还有什么比他更可信的证人呢？他了解这些情况，并如此记载了他们。\par

\SourceRecord{Translated by M. Atkinson and Archibald Robertson.；Nicene and Post-Nicene Fathers, Second Series；Vol. 4.；Edited by Philip Schaff and Henry Wace.；Buffalo, NY: Christian Literature Publishing Co.,；1892.；<http://www.newadvent.org/fathers/28151.htm>.}

% structure: p0001 source=h1 target=section reason=page-title

\section{\WorkTitle{亚略派历史，第二部分}}

\Signature{圣亚大纳削}\par

% structure: p0003 source=h2 target=subsection reason=relative-heading-rank; base=h2

\subsection{九、君士坦提乌斯治下的首次亚略派迫害}

然而，\Person{欧瑟比乌斯}及其同党看到他们的异端日渐式微，便写信到\Place{罗马}，并致函皇帝\Person{君士坦丁}和\Person{君士坦斯}，控告\Person{亚大纳削}。但当\Person{亚大纳削}所派遣的人驳倒了他们所写的陈述时，他们在皇帝面前蒙羞。\Place{罗马}主教\Person{儒利乌斯}也写信说，应在任何我们期望的地方召开一次会议，以便他们能提出他们所要控告的罪名，也能就他们所被指控的事自由辩护。他们派去的司铎看到自己在揭露真相，也请求这样做。于是，这群所作所为皆可疑的家伙，见自己在教会审判中不占上风，便投奔\Person{君士坦提乌斯}一人，从此向他哭诉，视他为异端的庇护者。他们说：\Quote{请宽免异端吧；您看，众人都离弃了我们，如今我们只剩下寥寥数人。请您开始迫害吧，因为我们连这仅剩的几人也舍弃了，我们落得孤苦无依。那些当年我们趁这些人被驱逐时拉拢过来的人，如今他们一回来便说服他们再度与我们作对。所以，请写信攻击他们所有人，再次派遣\Person{非拉格里乌斯}出任埃及总督，因为他能够为我们顺利地推行迫害，正如他先前已经试验过的那样，而且他身为背教者，更是如此。同时，请派遣\Person{额我略}前往\Place{亚历山大}任主教，因为他也能强化我们的异端。}\par

% structure: p0005 source=h2 target=subsection reason=relative-heading-rank; base=h2

\subsection{十、额我略的暴力闯入}

于是，\Person{君士坦提乌斯}立即写信，开始迫害所有人，并派遣\Person{非拉格里乌斯}出任总督，随行有一名叫\Person{阿尔萨奇乌斯}的宦官；他还派遣\Person{额我略}率领一支军队。随后发生了与先前一样的后果。他们纠集了一群牧人、羊倌以及镇上的其他放荡青年，手持刀剑棍棒，成群结队地攻击那座称为基里努斯教堂的教会；他们有的杀人，有的践踏，有的鞭打并投入监狱或流放。他们还拖拽许多妇女，公然将她们拖进法庭，侮辱她们，揪着她们的头发拖行。有些人被剥夺公权；有些人被夺去口粮，不为别的原因，只为诱使他们加入亚略派，并接纳皇帝派来的\Person{额我略}。\par

% structure: p0007 source=h2 target=subsection reason=relative-heading-rank; base=h2

\subsection{十一、东方人拒绝罗马会议}

然而，\Person{亚大纳削}在这些事发生之前，一听说他们的图谋，便启航前往\Place{罗马}，既知异端者的暴怒，也为了按先前决定召开会议。\Person{儒利乌斯}写信给他们，派遣司铎\Person{埃尔皮迪乌斯}和\Person{斐洛克塞努斯}，指定一个日期，让他们要么前来，要么自视为全然可疑之人。但是，当\Person{欧瑟比乌斯}及其同党一听说审判将是教会的审判，没有伯爵在场，门前也不驻扎士兵，而且诉讼不受皇家命令的约束（因为他们一向依赖这些来对抗主教，没有这些他们甚至不敢开口）；他们便如此惊恐，以至于拘留司铎们直到预定日期之后，并编造了一个不体面的借口，说他们因波斯人挑起的战争而不能来。但这并非他们拖延的真正原因，而是他们的良心恐惧。因为主教与战争有何相干？倘若他们因波斯人而不能来\Place{罗马}，尽管\Place{罗马}路途遥远且远隔重洋，那为何他们像狮子\BibleRef{伯前 5:8}一样，在东方的各处及靠近波斯人的地区四处巡行，寻找反对他们的人，以便诬告并放逐他们呢？\par

无论如何，当他们用这个不合理的借口打发走司铎们之后，他们彼此商议道：\Quote{既然我们在教会审判中无法占得上风，就让我们施展我们惯常的胆大妄为吧。}于是他们写信给\Person{非拉格里乌斯}，促使他不久后与\Person{额我略}一同前往埃及。随之，主教们遭到严厉鞭打，并被投入锁链。例如，主教兼精修圣人\Person{撒辣帕蒙}被他们驱逐流放；主教兼精修圣人\Person{波塔蒙}，在以往的迫害中已经失去一只眼睛，他们残酷地鞭打他的颈部，以至于在鞭打结束前他已看似死去。他被扔在一边，几乎过了几个钟头，经过细心照料和扇风，才苏醒过来，天主保全了他的性命；但没过多久，他便因鞭伤而死去，在基督内获得了第二次殉道的荣耀。除此之外，当\Person{额我略}与\InnerQuote{公爵}\Person{巴拉奇乌斯}一起坐着时，有多少隐修士被鞭打！多少主教受伤害！多少贞女被殴打！\par

% structure: p0010 source=h2 target=subsection reason=relative-heading-rank; base=h2

\subsection{十三、额我略在亚历山大的暴行}

此后，那个邪恶的\Person{额我略}号召所有人都要与他共融。但如果你要求他们共融，他们就不该受鞭打；如果你把他们当作恶人责打，为什么又把他们当作圣者要求共融呢？然而，他除了完成派遣他之人的计谋，建立异端之外，别无他图。因此，他愚昧地成了杀人犯和刽子手，恶毒、狡诈、亵渎；一言以蔽之，基督的仇敌。他如此残忍地迫害主教的姑母，以至于她死后他都不许她下葬。这本就是她的命运；她本会被弃尸不葬，若非那照看尸身的人把她当作自己的亲族抬走。即使在这等事上，他也显出亵渎的性情。再者，当寡妇和其他乞丐领到施舍时，他竟命令夺走已给他们的东西，并打碎他们盛油装酒的器皿，这样他不仅通过抢劫显示不虔，更在其行为中侮辱主；不久之后，他必将从主那里听到这些话：\Quote{你们既轻慢这些人，就是轻慢我。}\par

% structure: p0012 source=h2 target=subsection reason=relative-heading-rank; base=h2

\subsection{十四、额我略的亵渎与巴拉奇乌斯之死}

他还做了许多其他事，超乎语言所能描述，任何人听了都会觉得难以置信。他之所以如此行事，是因为他没有按照教会规则领受\OriginalTerm{祝圣}{ordination}，也没有按传统被召叫为主教；而是带着军事权威和排场，从宫廷被差遣而来，好像受托管理世俗政府的人。因此，他炫耀自己是总督们的朋友，胜过作主教和隐修士的朋友。所以，每当我们的父亲\Person{安当}从山中写信给他时，正如虔诚为罪人所憎恶，他也厌恶这位圣人的信件。但每当皇帝、将军或其他长官写信给他时，他就极其喜乐，就像\BibleBook{}{箴言}中所说的那些人，\OriginalTerm{圣言}{the Word}曾愤慨地论及他们说：\BibleQuote{祸哉，那些离弃正直之道、喜爱作恶、因恶人的乖僻而欢乐的人！} 他便重赏这些信件的使者；但有一次，安当写信给他，他竟指使\Person{巴拉基乌斯}公爵向那封信吐唾沫，并将它丢弃。然而，\OriginalTerm{天主的正义}{Divine Justice}并没有放过此事；因为不久之后，那公爵骑马前往第一站时，马转过头来，咬了他的大腿，把他摔下；三天之内他就死了。\par

\SourceRecord{Translated by M. Atkinson and Archibald Robertson.；Nicene and Post-Nicene Fathers, Second Series；Vol. 4.；Edited by Philip Schaff and Henry Wace.；Buffalo, NY: Christian Literature Publishing Co.,；1892.；<http://www.newadvent.org/fathers/28152.htm>.}

% structure: p0001 source=h1 target=section reason=page-title

\section{亚略派史，第三部分}

\Signature{\Emph{圣亚大纳削}}\par

% structure: p0003 source=h2 target=subsection reason=relative-heading-rank; base=h2

\subsection{十五、\Place{萨迪加}公会议上公教徒的恢复}

当他们对所有人采取类似手段时，在\Place{罗马}约有五十位主教集会，谴责\Person{尤西比乌斯}及其同党为可疑之人，不敢前来，并判定他们送来的书面陈述不足采信；而接纳了我们，欣然与我们共融。当这些事发生之际，关于\Place{罗马}公会议的召开、以及\Place{亚历山大里亚}及整个\Place{东方}教会遭受迫害的消息，传到了\Person{君士坦斯}皇帝的耳中。他致函其弟\Person{君士坦提乌斯}，两人随即决定召开公会议，以求解决争端，使受迫害者得以摆脱进一步的苦难，而迫害者无法再施暴行。于是，在\Place{萨迪加}城，由\Place{东方}与\Place{西方}共约一百七十位主教聚集；来自\Place{西方}的只是主教，以\Person{霍西乌斯}为他们的父执，但来自\Place{东方}的却带来了青年教师的辅佐和辩护律师，以及\Person{穆索尼安努斯}伯爵与卡斯托伦西斯的\Person{赫西基乌斯}；他们正是仗着此等人的声势，满心欣然地前来，以为一切又将由他们的权威所操控。因为藉着这些人，他们一向在所意图恐吓的人面前显得可畏，且对他们所选定的人肆意推行诡计。但当他们到达后，却发现案件将纯粹按教会事务处理，没有伯爵或军队的干预；当他们看见来自各教会各城的指控者，以及呈上的对他们不利的证据，当他们看见同来的可敬主教\Person{亚略}和\Person{阿斯特里乌斯}，竟离开他们而站到我们一边，并揭露他们的诡诈，指出他们的行径何等可疑，且因畏惧审判的结果而退缩，深恐向我们揭穿他们是虚伪的控告者，并使他们所派来的所谓控告人，暴露其所言全由他们授意，而他们乃是阴谋的策划人。觉察到这般情形，他们虽曾满怀热忱而来，以为我们必畏怯不敢与他们对质，但如今看见我们的从容，便将自己关闭在宫殿里（因他们的住处即在那里），然后彼此商议如下：\Quote{我们原是为一种结果而来，却看见了另一种结果；我们与伯爵们同行前来，审判却抛开他们进行。我们注定要被定罪了。你们都知道下达的命令。\Person{亚大纳削}和他的同道握有在\Place{马雷奥提斯}的记录，凭此他已得以澄清，而我们则蒙受耻辱。我们为何还要拖延？为何如此迟缓？让我们编造个借口离开吧，若留下必被定罪。宁可背负逃跑的羞耻，也胜过被揭露为诬告者的丑行。如果我们逃走，还能找到方法为我们的异端辩护；即使他们因我们逃亡而定我们的罪，我们仍有皇帝作靠山，皇帝不会容让民众将我们从教会里驱逐出去。}\par

% structure: p0005 source=h2 target=subsection reason=relative-heading-rank; base=h2

\subsection{十六、东方人士在\Place{萨迪加}的退出}

他们如此彼此商议之后，\Person{霍西乌斯}及所有其他主教一再向他们表示\Person{亚大纳削}及其同道之从容，说：\Quote{他们已准备好辩护，并保证证明你们是诬告者。} 又说：\Quote{你们若畏惧审判，又何以来此与我们相会？你们本不该来，或既已来了，就不该逃跑。} 当他们听到这话，越发惊恐，便求助一个较之在\Place{安提约基雅}所托的更为不堪的借口，即声称因皇帝写信告知他们他战胜波斯人的消息，故不得不离去。他们竟不惜差遣\Place{萨迪加}教会的长老\Person{尤斯塔提乌斯}去传递这一借口。然而即便如此，他们的逃跑并未如愿以偿；因为以伟大的\Person{霍西乌斯}为主席的圣公会议立时致函向他们明确表示，说：\Quote{你们或是前来，就所提出的控诉，及你们向他人所捏告的答复，否则须知道，公会议将判你们有罪，并宣布\Person{亚大纳削}及其同道清白无罪，毫无过失。} 但他们因良心的惊惧，非但不顺应信函的提议，反而更加迅速逃跑；当他们看见那些曾受他们侵害的人时，甚至不敢转过脸去听他们的话，反而跑得更快。\par

% structure: p0007 source=h2 target=subsection reason=relative-heading-rank; base=h2

\subsection{十七、\Place{萨迪加}公会议的决议}

就在如此可耻而不堪的景象下，他们逃跑了。圣公会议由超过三十五个省区聚集而成，洞悉亚略派的恶意，遂允许\Person{亚大纳削}及其同道，针对他们所提的指控进行答辩，并陈述所遭受的苦难。当他们如此辩护完毕，如我们前文所述，公会议赞许并高度钦佩他们的操守，欣然接纳他们的共融，并发函至各处，至每一教区，尤其致函\Place{亚历山大里亚}、\Place{埃及}及\Place{利比亚}地区，宣告\Person{亚大纳削}及其友人清白无罪，毫无过失，而对手则是诽谤者、作恶者，乃至不配称为基督徒之人。于是，他们在平安中被遣散去；但公会议罢黜了\Person{斯特凡}、\Person{梅诺凡特}、\Person{阿卡基乌}、\Place{劳狄刻雅}的\Person{乔治}、\Person{乌尔萨西}、\Person{瓦伦特}、\Person{提奥多尔}及\Person{那耳索斯}。对于那位被皇帝派往\Place{亚历山大里亚}的\Person{额我略}，公会议发表了一篇宣告，称他从未被祝圣为主教，且不应被称为基督徒。因此，公会议声明他所自称授予的圣秩无效，并因其新奇而不合法之性质，下令在教会中甚至不可提及。就这样，\Person{亚大纳削}及其友人在平安中被遣散（这些相关信函因其篇幅冗长而附于书末），公会议便解散了。\par

% structure: p0009 source=h2 target=subsection reason=relative-heading-rank; base=h2

\subsection{\Place{撒狄卡}会议后的亚略派迫害}

但那些本应保持沉默的被罢免者，以及那些在如此可耻的逃跑后分离的人，他们的行为如此恶劣，以至于他们之前的所作所为与此相比显得微不足道。因为当\Place{哈德良堡}的民众不愿与他们共融时——视他们为逃避会议并已被判有罪的人——他们便向皇帝\Person{君士坦提乌斯}申诉，并成功使那里的兵器制造厂的十名平信徒被斩首，当时再次担任伯爵的\Person{菲拉格留斯}也在这件事上协助了他们的计划。我们在经过时曾见过这些人的坟墓，它们位于城外。然后，仿佛他们因自己的逃跑而免于被定罪为诬告者便大获成功，他们说服皇帝下令，凡他们所愿的都要执行。这样，他们使两位司铎和三位执事从\Place{亚历山大里亚}被流放到\Place{亚美尼亚}。至于\Person{亚略}和\Person{阿斯忒里乌斯}，一位是\Place{巴勒斯坦}的\Place{培特辣}的主教，另一位是\Place{阿拉伯}的主教，因脱离他们的党派，他们不仅将他们流放到\Place{上利比亚}，还使他们遭受侮辱。\par

% structure: p0011 source=h2 target=subsection reason=relative-heading-rank; base=h2

\subsection{对\Place{亚历山大里亚}人的暴虐措施}

至于\Place{哈德良堡}的主教\Person{路西乌斯}，当他们发现他大胆发言反对他们，揭露他们的不虔时，便像以前一样，使他颈部和双手被铁链锁住，然后将他驱逐流放，并在那里死去，这是他们所知道的。他们还罢免了主教\Person{狄奥多罗斯}；而对于\Place{色雷斯}两位善良正统的主教，\Place{艾诺斯}的\Person{奥林皮乌斯}和\Place{图拉真堡}的\Person{狄奥多鲁斯}，当他们察觉到他们对异端的憎恨时，便提出了虚假的指控。这件事\Person{欧瑟比乌斯}及其同伙最先发难，皇帝\Person{君士坦提乌斯}就此写信；随后这些人又重提指控。信中的大意是，他们不仅被逐出他们的城市和教会，而且无论在哪里被发现，都应处死刑。无论这种行为多么令人惊讶，这完全符合他们的原则；因为他们受了\Person{欧瑟比乌斯}及其同伙的这些做法的训练，又继承了他们的不虔和恶毒原则，便想在\Place{亚历山大里亚}显出自己的威风，如同他们的父辈曾在\Place{色雷斯}所做的那样。他们下令监视港口和城门，以免那些被流放的人利用会议所给予的许可返回自己的教会。他们还下令给\Place{亚历山大里亚}的地方长官，关于\Person{亚大纳削}及其中指名的某些司铎，规定如果这位主教或其他任何人被发现来到该城或其边境，长官有权将发现的人斩首。因此，这种新的犹太式异端不仅否认主，而且学会了杀人。\par

% structure: p0013 source=h2 target=subsection reason=relative-heading-rank; base=h2

\subsection{在\Place{安提约基雅}针对公教使节的阴谋}

然而，即使在这之后他们仍未停止；正如他们的异端之父如同咆哮的狮子巡游，寻找可吞噬的人，这些人获得使用驿马的权利到处奔走，无论何时发现有人指责他们逃跑、憎恶亚略异端，他们便鞭打他们、给他们戴上锁链，将他们驱逐出境；他们变得如此可怕，以致许多人不得不伪装，许多人逃往旷野，甚至不愿与他们有任何交往。这就是他们疯狂驱使他们在逃跑后所犯的暴行。此外，他们还施行了另一桩令人发指的恶行，这确实符合他们异端的特性，却是我们前所未闻、即便在更放荡的外邦人中也不大可能再度发生的，更不用说在基督徒中了。神圣的会议曾派遣两位主教作为使节，即\Place{卡普阿}的\Person{文森提乌斯}（这是\Place{坎帕尼亚}的都城）和\Place{阿格里皮纳}的\Person{欧弗辣提斯}（这是\Place{上高卢}的都城），为要获得皇帝对会议决议的同意，使主教们得以返回自己的教会，因为他是他们被驱逐的始作俑者。至为虔诚的\Person{君士坦斯}也曾写信给他的兄弟，支持主教们的事业。但这些无恶不作的人，当他们看到两位使节在\Place{安提约基雅}时，便共谋策划，\Person{斯特凡努斯}自告奋勇独自去执行，认为自己是达此目的的合适工具。于是，他们雇了一个普通的妓女，甚至在最神圣的复活节期间，让她赤身，夜间领她进入主教\Person{欧弗辣提斯}的房间。那妓女原以为是一位年轻人派人来请她，一开始欣然跟随他们；但当他们把她推进去，她看到那人睡熟了，对所发生的事全然不知，随后她看清了那人的容貌，见到是一位老人的面孔和主教的装束时，她立刻大声喊叫，声称自己遭到了暴力。他们要求她噤声，并对主教作出虚假指控；这样到了白天，此事便传开，全城的人都聚集起来，从皇宫来的人大为骚动，对所传播的消息感到惊讶，并要求不可对此沉默不语。因此进行了调查，她的主人提供了有关那些来接妓女之人的信息，而这些人又告发了\Person{斯特凡努斯}；因为他们正是他的圣职人员。于是，\Person{斯特凡努斯}被罢免，宦官\Person{莱翁提乌斯}被任命接替他的位置，只为亚略异端不缺乏一个支持者。\par

% structure: p0015 source=h2 target=subsection reason=relative-heading-rank; base=h2

\subsection{\Person{君士坦提乌斯}改变主意}

如今，\Person{君士坦提乌斯}皇帝感到一些良心不安，便恢复常态；他从他们对\Person{优弗拉特斯}的行为断定，他们对其他人的攻击也是同类的，于是下令，将先前从\Place{亚历山大里亚}放逐到\Place{亚美尼亚}的司铎们和执事们立即释放。他还公开写信给\Place{亚历山大里亚}，命令不得再迫害\Person{亚大纳削}的友人们，无论是圣职人员还是信友。大约十个月后，当\Person{格里高利}死去时，他便以一切尊荣之礼派人去请\Person{亚大纳削}，给他写了不下三次非常友好的信，勉励他勇敢前来。他还派遣一位司铎和一位执事，为使他更受鼓舞而返回；因为他认为，由于先前发生之事带来的惊恐，我不愿返回。此外，他还写信给他的兄弟\Person{君士坦斯}，请他也劝我返回。并且他断言，他已期待\Person{亚大纳削}整整一年，他不允许任何变动或任何祝圣举行，因为他正为作为它们主教的\Person{亚大纳削}保留着各教会。\par

% structure: p0017 source=h2 target=subsection reason=relative-heading-rank; base=h2

\subsection{二十二、\Person{亚大纳削}拜访\Person{君士坦提乌斯}}

因此，当他以这种语气写信，并借着许多人鼓励他（因为他让\Person{波勒米乌斯}、\Person{达提阿努斯}、\Person{巴狄翁}、\Person{塔拉苏斯}、\Person{陶鲁斯}和\Person{弗洛伦提乌斯}——这些\Person{亚大纳削}最能信赖的侍臣——也都写信），\Person{亚大纳削}便将整个事情交托于那位触动\Person{君士坦提乌斯}良心行事的天主，与朋友们一起来到他那里；皇帝给了他友善的接见，打发他返回自己的家乡和教会，同时给各地长官写信，说先前他曾命令把守道路，现在应准其自由通行。随后，当主教抱怨自己曾受的苦，以及皇帝所写攻击他的那些信函，并恳求不要让他的仇敌在他离开后再提起那些诬告，说：\Quote{若您愿意，请传唤这些人；因为就我们而言，他们可以自由站出，我们要揭露他们的行为。} 他不愿这样做，却命令将以前诽谤他的文字全部销毁抹除，断言他永不再听信此类控诉，他的决定是坚定而不可更改的。他并非仅仅这样说，更是呼求天主作证，用誓言封缄了他的话。于是，他又用许多别的话鼓励他，要他满怀信心，便向主教们及各官员发出以下信函。\par

23. \Person{君士坦提乌斯}·奥古斯都，大帝，征服者，致天主教会的主教及圣职人员。\par

最可敬的\Person{亚大纳削}并未被天主的\OriginalTerm{恩宠}{grace}所遗弃，等等。\par

\Emph{另一封信。}\par

\Person{君士坦提乌斯}致\Place{亚历山大里亚}人民。\par

我们既在各方面祝愿你们安康，等等。\par

\Emph{另一封信。}\par

\Person{君士坦提乌斯}·奥古斯都，征服者，致\Place{埃及}总督\Person{涅斯托里乌斯}。\par

众所周知，我们曾下达过一道命令，且发现某些文件有损最可敬的\Person{亚大纳削}主教之名誉；这些文件存在于贤明阁下的\OriginalTerm{指令簿}{Order-book}中。因此，我们期望审慎贤明的阁下（我们对此已有确证）遵从此令，将存于\OriginalTerm{指令簿}{Order-book}中、涉及前述人物的所有信函，转呈于我们的朝廷。\par

24. 以下是他于蒙福的\Person{君士坦斯}去世后所写的信函。该信原以拉丁文写就，此处译成希腊文。\par

\Person{君士坦提乌斯}·奥古斯都，征服者，致\Person{亚大纳削}。\par

你的审慎并非不知，我恒常祈祷，愿我已故的兄弟\Person{君士坦斯}在他一切事业中凡事顺遂；因此你的智慧当能想见，当我得知他被极其渎神之手夺去生命时，我何其悲痛。如今，正当这真正哀恸的时刻，有些人正试图惊扰你，因此我认为宜修书一封致你的坚毅，劝勉你，身为一位主教，当教导人民那些关乎神圣宗教之事，并如你所常行的那样，同他们一起以祈祷度日，而不要轻信任何虚妄的传闻。因为我们已定意，你当依照我们的愿望，在你自己的地方继续履行主教之职。愿\OriginalTerm{天主眷顾}{Providence}保全你，至爱的父亲，多年之久。\par

% structure: p0030 source=h2 target=subsection reason=relative-heading-rank; base=h2

\subsection{二十五、\Person{亚大纳削}从第二次流放归来}

在此情形下，当他们最终告别、启程之后，他的朋友们为见到友人而欢欣；而对方一党，有些人见到他便感到羞愧；另一些人不敢露面，藏了起来；还有些人悔恨自己曾写下的反对主教的文字。于是，\Place{巴勒斯坦}的所有主教，除两三个品行可疑者外，都欣然接纳了\Person{亚大纳削}，与他共融，甚至写信为自己辩解，声称他们先前所写的，并非出于己愿，而是迫于压力。至于\Place{埃及}和\Place{利比亚}诸省的主教，以及这些地区和\Place{亚历山大里亚}的平信徒，我无需赘述。所有人都聚集而来，怀着难以言表的喜悦，因为他们不仅出乎意料地见到朋友们活着，而且摆脱了那些如暴君和狂犬般对待他们的异端者。因此，他们的喜乐极大，会众在聚会中彼此勉励修德。多少先前准备结婚的少女，如今为基督守贞！多少青年看见他人的榜样，拥抱了隐修生活！多少父亲劝导子女，又有多少被子女催促，不要阻挠基督徒的克修！多少妻子说服丈夫，又有多少被丈夫说服，如宗徒所言\BibleRef{格前 7:5}，专心祈祷！多少寡妇孤儿，先前饿着赤身，如今因着民众极大的热忱，不再饥饿，且有衣穿！简言之，他们在修德上的争竞如此之大，以致你会以为每个家庭、每座房屋都是一所教会，因其成员的善行和向天主的祈祷。在教会中，有着深邃而奇妙的和平，各地主教纷纷致信，并依照惯例，领受\Person{亚大纳削}的和平信函。\par

% structure: p0032 source=h2 target=subsection reason=relative-heading-rank; base=h2

\subsection{二十六、\Person{乌尔撒奇}与\Person{瓦伦斯}的悔改}

此外，\Person{乌尔撒奇}与\Person{瓦伦斯}仿佛良心受到鞭策，回心转意，他们虽然未收到\Person{亚大纳削}的任何信息，却主动写信给他，措辞友善和睦。他们又前往\Place{罗马}，懊悔并承认先前对主教的种种指控和说法，全是虚谎与诽谤。他们不仅自愿如此行，还绝罚了\OriginalTerm{亚略异端}{Arian heresy}，并呈上书面悔过声明，以拉丁文致信\Person{犹利乌斯}主教，后译为希腊文。拉丁文副本由\Place{特里尔}的主教\Person{保禄}寄送给我们。\par

\Emph{拉丁文译文。}\par

\Person{乌尔撒奇}与\Person{瓦伦斯}致我主至福教宗\Person{犹利乌斯}。\par

众所周知，我们云云。\par

\Emph{拉丁文译文。}\par

主教\Person{乌尔撒奇}与\Person{瓦伦斯}致我主暨兄弟\Person{亚大纳削}主教。\par

适逢递呈之机，云云。\par

写完这些之后，他们还在\Person{亚大纳削}的司铎\Person{伯多禄}与\Person{依勒内}，以及路过该地的平信徒\Person{阿摩尼}向他们呈递的和平信函上署名；尽管\Person{亚大纳削}甚至连通过这些人也未曾给他们带过任何口信。\par

% structure: p0041 source=h2 target=subsection reason=relative-heading-rank; base=h2

\subsection{二十七、\Person{亚大纳削}的凯旋}

目睹这一切，以及教会中盛行的莫大和平，谁不充满赞叹？看到如此众多主教的和谐，谁不欢喜？见到集会中民众的喜悦，谁不光荣主？多少敌人悔改了！多少先前诬告他的人为自己辩解！多少曾经恨他的人，如今对他示好！多少写信反对他的人，收回了他们的陈述！许多并非因选择、而是因被迫而站在\OriginalTerm{亚略派}{Arians}一边的人，夜间前来，为自己辩白。他们绝罚了异端，恳求他的宽恕，因为虽因这些人的阴谋与诽谤，他们在肉身上看似站在他们一边，但内心始终与\Person{亚大纳削}相通，一直与他同在。\Emph{请相信我，这是真的。}\par

\SourceRecord{Translated by M. Atkinson and Archibald Robertson.；Nicene and Post-Nicene Fathers, Second Series；Vol. 4.；Edited by Philip Schaff and Henry Wace.；Buffalo, NY: Christian Literature Publishing Co.,；1892.；<http://www.newadvent.org/fathers/28153.htm>.}

% structure: p0001 source=h1 target=section reason=page-title

\section{亚略派历史 第四部分}

\\Signature{圣亚大纳修著}\par

% structure: p0003 source=h2 target=subsection reason=relative-heading-rank; base=h2

\subsection{君士坦提乌斯治下的第二次亚略派迫害}

但是\\Person{优西比乌}及其党徒的意见与不虔敬的继承人，即太监\\Person{莱翁提乌斯}——他甚至作为平信徒也不配留在共融内，因为他自残，以便从此可以自由地与一位名叫\\Person{优斯托利乌姆}的女子同床，她实际上作他的妻子，却被称为童贞女——还有被公会议罢免的\\Person{乔治}、\\Person{阿卡西乌斯}、\\Person{狄奥多鲁斯}和\\Person{纳尔基索斯}，他们听说并看见这些事，就极其羞愧。当他们察觉到\\Person{亚大纳修}与众主教之间所存在的同心合意与平安时（这些主教超过四百位，来自伟大的\\Place{罗马}、整个\\Place{意大利}、\\Place{加拉布里亚}、\\Place{阿普利亚}、\\Place{坎帕尼亚}、\\Place{布鲁提亚}、\\Place{西西里}、\\Place{撒丁尼亚}、\\Place{科西嘉}以及整个\\Place{阿非利加}；来自\\Place{高卢}、\\Place{不列颠}和\\Place{西班牙}的，与伟大的宣信者\\Person{霍西乌斯}同来；还有来自\\Place{帕诺尼亚}、\\Place{诺里库姆}、\\Place{西斯奇亚}、\\Place{达尔马提亚}、\\Place{达尔达尼亚}、\\Place{达西亚}、\\Place{美西亚}、\\Place{马其顿}、\\Place{帖撒利}、整个\\Place{亚该亚}，以及\\Place{克里特}、\\Place{塞浦路斯}、\\Place{吕基亚}的，还有\\Place{巴勒斯坦}、\\Place{伊索里亚}、\\Place{埃及}、\\Place{底比斯}、整个\\Place{利比亚}和\\Place{彭塔波利斯}的大部分主教），当他们察觉到这些事——我这样说——他们就被嫉妒与恐惧所充满：嫉妒，是因为如此众多的人聚集共融；恐惧，是因为唯恐那些被他们诱入陷阱的人，会被如此众多人的同心合意所争取过去，从此他们的异端便大遭挫败，声名扫地，在各处都被禁止。\par

% structure: p0005 source=h2 target=subsection reason=relative-heading-rank; base=h2

\subsection{乌尔撒西乌斯与瓦伦斯的重新堕落}

首先他们说服\\Person{乌尔撒西乌斯}、\\Person{瓦伦斯}及其党徒再次改变立场，就像\\BibleQuote{狗转过身来吃自己所吐的，又如猪洗净了又回到泥里去打滚}，回到他们先前不虔的污秽中\\BibleRef{伯后 2:22}；并且他们为自己的变节找借口，说是出于对至虔敬的\\Person{君士坦斯}的畏惧。然而，即使确有畏惧的原因，如果他们对自己所行的事有信心，也不该背叛自己的朋友。现今既无畏惧的理由，而且他们是在说谎，难道不是该受完全的定罪吗？因为当时没有士兵在场，没有像他们现在这样派遣侍卫或书记官，皇帝也不在场，也没有任何人邀请他们，他们就写下了悔改书。相反，他们自愿上到\\Place{罗马}，在教会里自己主动悔改并记录下来，那里没有外来的恐惧，唯一的恐惧是对天主的敬畏，在那里人人都有良心的自由。他们虽然第二次成了亚略派，却又为自己的行为编造出如此不成体统的借口，仍全然不知羞耻。\par

% structure: p0007 source=h2 target=subsection reason=relative-heading-rank; base=h2

\subsection{君士坦提乌斯再次改变立场}

接下来他们成群去见皇帝\\Person{君士坦提乌斯}，恳求他说：\\Quote{当初我们初次向你请求时，我们并未得到信任；因为我们曾对你说过，在你召见\\Person{亚大纳修}之时，通过邀请他前来，你正在驱逐我们的异端。因为他从一开始就反对它，并且从未停止诅咒它。他已经写信到世界各地反对我们，大多数人已经与他共融；甚至那些看似站在我们这边的人，有些已被他争取过去，其余的人也可能如此。我们便被孤立了，以致我们担心，唯恐我们的异端的真面目被人知晓，从此以后，我们和你都会被冠以异端者的名号。若果真要如此，你须留意，使我们不致被归于摩尼教徒之列。因此，重新开始迫害，支持异端吧，因为它认你作它的君王。}这就是他们邪恶的说辞。皇帝在匆忙进军对抗\\Person{马格嫩提乌斯}而经过该地区时，目睹了众主教与\\Person{亚大纳修}的共融，就如火灼一般，突然改变了主意，不再记得自己的誓言，同样忘记了他所写过的书信，也不顾对他兄弟应尽的义务。因为在他写给他兄弟的信中，以及在他与\\Person{亚大纳修}的会面中，他都起誓说，他不会做出违背人民意愿以及众主教心意的事。但对不虔的热忱使他立刻忘掉了这一切。然而，我们不必惊奇，\\Person{君士坦提乌斯}在写了那么多信、发了那么多誓后竟然改变了主意，因为我们记得从前的\\Person{法老}，\\Place{埃及}的暴君，在屡次承诺并因而获得惩罚的暂缓后，也同样改变，直到最后他与他的同伙一同灭亡\\BibleRef{出 14:28}。\par

% structure: p0009 source=h2 target=subsection reason=relative-heading-rank; base=h2

\subsection{君士坦提乌斯开始迫害}

他于是强迫各城的民众改变党派；到达\Place{阿尔勒}和\Place{米兰}后，他完全依照异端者的图谋和建议行事；或者更确切地说，他们自己采取行动，从他那里获得权力后，猛烈地攻击每一个人。信件和命令立刻被送到此地的长官那里，要求今后将粮食从\Person{亚大纳削}那里收回，交给那些拥护亚略派教义的人，并且无论谁都可以随意侮辱那些与他共融的人；地方官若不与亚略派共融，便会受到威胁。这些事不过是随后在公爵\Person{叙利亚努斯}指挥下所发生之事的前奏。命令也被送到了更偏远的地方，御前公证人被派往每一个城市，还有宫内官，带着对主教和地方官的威胁，指示地方官催促主教，并告知主教，他们要不在反对\Person{亚大纳削}的文件上签名，并与亚略派共融，要不就自己承受流放的惩罚，而与他们站在一起的民众则要明白，锁链、凌辱、鞭打和丧失财产，将是他们的下场。这些命令并未被忽视，因为委员们的随行中有\Person{乌尔撒西乌斯}和\Person{瓦伦斯}的神职人员，以激发他们的热忱，并在地方官失职时向皇帝报告。其他的异端，作为他们自己的姐妹，他们容许其亵渎主，唯独对基督徒合谋攻击，无法容忍听到关于基督的正统言论。结果，按照圣经的话，有多少主教被带到官长和君王面前（\BibleRef{谷 13:9}），并从地方官那里得到这样的判决：\Quote{签署吧，否则离开你们的教会，因为皇帝已下令废黜你们！}在各城里有多少人受到粗暴对待，免得他们控诉那些主教的支持者！更有信件被送至城市当局，并以罚款相威胁，如果他们没有强迫各自城里的主教签署的话。总之，每个地方和每个城市都充满了恐惧和混乱，主教们被拖去受审，地方官则目睹民众的悲恸与哀叹。\par

% structure: p0011 source=h2 target=subsection reason=relative-heading-rank; base=h2

\subsection{君士坦提乌斯的迫害}

这就是宫内委员们的所作所为；另一方面，那些可钦佩的人，对自己已获得的庇护深信不疑，表现出极大的热忱，致使一些主教被传唤到皇帝面前，同时他们又写信迫害另一些主教，捏造指控；为的是使前者慑于\Person{君士坦提乌斯}的威临，后者则因恐惧委员们和这些虚假指控中所含的威胁，而被迫放弃他们正统而虔诚的见解。皇帝就是用这种方式，部分以威胁，部分以许诺，迫使如此众多的主教宣称：\Quote{我们不再与\Person{亚大纳削}共融。}因为那些前来觐见的人，既未蒙召见他，也不得片刻松懈，甚至连走出住所都不行，直到他们要么签署，要么拒绝并因此遭受放逐。他这样做，是因为他看到这异端为众人所憎恶。正因如此，他尤其强迫许多人将自己的名字加入那少数亚略派的行列，他热切渴望集结大批人名，既出于对这位主教的嫉妒，也是为了炫耀他所庇护的亚略派的不虔敬；他以为他能像影响人的思想一样轻易地改变真理。他不知道，也从未读过，撒杜塞人和黑落德党人，联合了法利塞人，仍不能模糊真理；相反，真理因此日复一日更加光辉，当他们喊着\Quote{我们除了\Person{凯撒}，没有君王}，并获得了\Person{比拉多}有利于他们的判决时，却仍被遗弃，在羞辱中等待，他们看到自己的庇护者临近死亡时，料想自己不久也将像鹧鸪一样失去一切。\par

% structure: p0013 source=h2 target=subsection reason=relative-heading-rank; base=h2

\subsection{迫害来自魔鬼}

如果主教们仅仅因为惧怕这些事就改变自己的意见，这完全不恰当，那么用暴力强迫不愿服从的人，就更加不恰当了，而且这不是那些对自己信仰有信心的人所当作的。这样，当魔鬼没有真理支持时，就用斧头和锤子攻击并砸开迎接他之人的大门。但我们的救主如此温柔，祂这样教导说：\BibleQuote{若有人愿意跟随我}，以及：\BibleQuote{谁若愿意作我的门徒}\BibleRef{玛 16:24}；祂来到每个人面前，并不强逼他们，反而敲门说：\BibleQuote{请给我开门，我的妹妹，我的爱卿}\BibleRef{歌 5:2}；如果他们给祂开门，祂就进去，但如果他们迟延不肯，祂就离开他们。因为真理不是靠刀剑、标枪，也不是靠士兵来传布的；而是靠着劝勉和规劝。但是，在畏惧皇帝的气氛笼罩下，哪有什么劝勉可言？当违抗他们的人最终遭到流放和死亡时，又哪有什么规劝？就连\Person{达味}，虽然自己是国王，又将敌人掌握在手中，当士兵们想要杀死他的敌人时，他并没有运用权威去制止他们，反而正如圣经所说，\Person{达味}用道理说服了手下的人，不许他们起来杀死\Person{撒乌耳}。\BibleRef{撒上 26:9} 但他，由于没有理性的论据，便用权力强迫所有人，为要显明给众人看：他们的智慧不是来源于天主，而纯粹是属人的，并且那些偏向\OriginalTerm{亚略派学说}{Arian doctrines}的人，实在是以\Person{凯撒}为他们的君王，而非以天主为君王；因为正是藉着\Person{凯撒}，这些基督的仇敌才能为所欲为。但是，正当他们以为藉着\Person{凯撒}向许多人推行他们的计谋时，他们却不知道，自己正使许多人成为\OriginalTerm{宣信者}{confessors}，其中就有那些最近作出如此光荣宣认的虔诚之士、卓越的主教们：\Place{高卢}都会\Place{特里尔}的主教\Person{保利努斯}、\Place{撒丁岛}都会主教\Person{路西法}、\Place{意大利}\Place{韦尔切利}的\Person{欧瑟伯}，以及\Place{意大利}都会\Place{米兰}的\Person{狄奥尼修斯}。皇帝把这些人召到面前，命令他们签字反对\Person{亚大纳削}，并与异端者共融；当他们被这新奇的做法惊得目瞪口呆，说并无这样的教会法典时，皇帝立即说：\Quote{凡我所意愿的，就当视为法典；\InnerQuote{叙利亚的主教们}让我这样说话。那么，你们要么服从，要么就去流放。}\par

% structure: p0015 source=h2 target=subsection reason=relative-heading-rank; base=h2

\subsection{西方主教们的流放传播了真理的知识}

主教们听到这话，大为惊愕，他们向天主伸开双手，鼓起极大的勇气向他直言，教导他说，国度不是他的，而是天主的，是天主赐给他的，他们还劝他畏惧天主，免得天主忽然将这国度从他夺去。他们用审判之日警戒他，警告他不要破坏教会的秩序，不要将罗马的国家权力与教会的宪制相混，不要将\OriginalTerm{亚略异端}{Arian heresy}引入天主的教会。但是他不肯听从他们，也不许他们多言，反而更加威胁他们，拔出剑来指着他们，并下令将其中一些人处死；虽然后来，就像\Person{法郎}一样，他又后悔了。因此，这些圣洁的人抖掉脚上的尘土，仰望天主，既不畏惧皇帝的威吓，也不在他拔出的剑前背叛他们所拥护的事业；而将流放视为自己牧职当尽之职分来接受。他们在沿途所经的每一处、每一座城，虽然身带锁链，却仍宣讲福音，宣认正统信仰，谴责\OriginalTerm{亚略异端}{Arian heresy}，并痛斥\Person{乌尔撒奇}和\Person{瓦伦斯}的变节。然而，这与他们仇敌的意图恰恰相反；因为这些被流放者所涉足之地越远，就越激起人们对施暴者的痛恨，而这些人的流离漂泊无非是在宣告施暴者的不虔敬。因为凡看见他们沿路经过的人，谁不深深敬仰这些\OriginalTerm{宣信者}{confessors}，而对那些人弃绝、憎恶，称他们不但是不虔敬的人，而且是刽子手、杀人犯，说是基督徒，毋宁说什么都不是？\par

\SourceRecord{Translated by M. Atkinson and Archibald Robertson.；Nicene and Post-Nicene Fathers, Second Series；Vol. 4.；Edited by Philip Schaff and Henry Wace.；Buffalo, NY: Christian Literature Publishing Co.,；1892.；<http://www.newadvent.org/fathers/28154.htm>.}

% structure: p0001 source=h1 target=section reason=page-title

\section{亚略派史，第五部分}

\Signature{圣亚大纳修 著}\par

% structure: p0003 source=h2 target=subsection reason=relative-heading-rank; base=h2

\subsection{三十五、\Person{利伯里乌}的迫害与失足}

35. 本来，如果\Person{君士坦提乌斯}从一开始就根本未曾与这异端有任何牵连，那倒更好；或者既已牵连，如果他不曾如此屈服于那些不虔之人，那也会更好；或者既已屈服，如果他只与他们一同到此为止，那么审判就只为这些暴行而临到他们众人了。可是，看来他们如同疯人一般，既已自陷于邪恶的捆绑中，便是在为自己招来更严厉的审判。因此，他们从一开始就连罗马主教\Person{利伯里乌}也没有放过，反而将他们的狂怒伸展到了那些地区；他们不尊重他的主教职位，因为那是一个\OriginalTerm{宗徒的宝座}{Apostolical throne}；他们对\Place{罗马}毫不敬畏，因为她是\OriginalTerm{罗马尼亚的都会}{Metropolis of Romania}；他们不记得，从前在他们的信函中，他们曾称罗马的主教们为\OriginalTerm{宗徒之人}{Apostolical men}。但既将一切混淆不分，他们便立刻忘却一切，只在乎表现出对不虔的热心。当他们察觉他是一位正统之人，且憎恶\OriginalTerm{亚略异端}{Arian heresy}，并竭力劝勉众人弃绝并远离这异端时，那些不虔之人便私下如此议论说：\Quote{我们若能说服\Person{利伯里乌}，我们便很快能胜过众人了。}于是，他们在皇帝面前诬告他；皇帝指望借\Person{利伯里乌}轻易将众人拉拢到他这边，便写信给他，并派了一个名叫\Person{犹西比乌}的太监带了信函与礼物前往，要用礼物讨好他，并拿信函威吓他。那太监于是到了\Place{罗马}，首先向\Person{利伯里乌}提议，要他签署反对\Person{亚大纳修}的文书，并与\OriginalTerm{亚略派}{Arians}相通共融，说道：\Quote{皇帝希望如此，并命令你这样做。}然后，他拿出礼物来，拉着他的手，又恳求他说：\Quote{顺从皇帝吧，收下这些礼物。}\par

% structure: p0005 source=h2 target=subsection reason=relative-heading-rank; base=h2

\subsection{三十六、太监\Person{犹西比乌}试图说服\Person{利伯里乌}而无效}

但主教力图说服他，这样对他推论说：\Quote{我怎能做这反对\Person{亚大纳修}的事呢？我们怎能定一个人的罪，这个人，不单一个公会议，而是由世界各地聚集的第二个公会议，已经公正地宣判他无罪，而且\Place{罗马}人的教会已平安地将他遣回？如果我们在他缺席时拒绝一个我们曾欣然欢迎其来到我们中间并接纳他共融的人，谁会赞成我们的行为呢？这并非教会的\OriginalTerm{法典}{Canon}；我们也没有从\OriginalTerm{教父们}{Fathers}领受过这样的传统，他们则是从伟大而蒙福的宗徒\Person{伯多禄}领受的。但如果皇帝真的关心教会的平安，如果他要求我们关于\Person{亚大纳修}的信函要予以撤销，那么，他们反对他及其他所有人的程序也应一并撤销；然后，应在远离宫廷之处召开一个\OriginalTerm{教会公会议}{Ecclesiastical Council}，皇帝不得出席，也不得允许任何\OriginalTerm{伯爵}{Count}或长官进来威吓我们，那里只应有对天主的敬畏和\OriginalTerm{宗徒的规范}{Apostolical rule}施行统治；这样，首先，教会的信仰才能稳固，如同教父们在\OriginalTerm{尼西亚大公会议}{Council of Nicæa}中所规定的那样，并且支持\OriginalTerm{亚略教义}{Arian doctrines}的人应被逐出，他们的异端应受\OriginalTerm{绝罚}{anathema}。然后，在那之后，再对\Person{亚大纳修}及其他任何人所遭受的指控，以及对方所受指控的事项进行调查，让有罪者被逐出，无辜者得到鼓励与支持。因为坚持不虔信条的人不可能被接纳为公会议的成员；而且在教义的问题之前，先调查行为的问题是不合适的；一切信仰观点上的分歧应首先被铲除，然后再调查行为的问题。我们的主\Person{耶稣基督}没有医治那些受苦难的人，直到他们表明并宣认他们对祂怀有何等信心。这些事，我们是从教父们领受的；将这些事报告给皇帝吧，因为这些事既对他有益，也能建树教会。但不要听从\Person{乌尔撒奇}和\Person{瓦伦提}，因为他们已撤回先前的声明，他们现在所说的，不可信靠。}\par

% structure: p0007 source=h2 target=subsection reason=relative-heading-rank; base=h2

\subsection{三十七、\Person{利伯里乌}拒绝皇帝的礼物}

这些就是主教\Person{利伯里乌}所说的话。那太监感到恼怒，倒不是因为他不肯签署，而是因为他发现他乃是异端的仇敌；他忘记自己是站在一位主教面前，在严厉地威胁了他之后，便带着礼物离去了；随后他犯下一件可耻的事，这事与基督徒身份不符，且对太监而言也过于狂妄。他效法\Person{撒乌耳}的过犯，就到宗徒\Person{伯多禄}的殉道纪念所，在那里献上礼物。但\Person{利伯里乌}得知此事后，对看守那地方的人非常生气，因为他没有阻止他，并把那些礼物当作非法的祭物扔出去，这更加激怒了那个被阉割之人对他的怒气。因此，他更是激怒皇帝来反对他，说：\Quote{我们所关心的事，已不再是取得\Person{利伯里乌}的签署，而是他如此坚决地反对异端，竟指名绝罚\OriginalTerm{亚略派}{Arians}。}他还煽动其他太监说同样的话；因为环绕\Person{君士坦提乌斯}身边的那些人，确切地说全部人数，都是太监，这些人完全垄断了他那里的影响力，没有他们，在那里任何事都办不成。皇帝于是写信到\Place{罗马}，再次派了\OriginalTerm{御前侍卫}{Palatines}、\OriginalTerm{书记官}{Notaries}和\OriginalTerm{伯爵们}{Counts}，带信给\OriginalTerm{总督}{Prefect}，致使他们要么用计谋将\Person{利伯里乌}诱离\Place{罗马}，把他送到宫廷他那里，要么以暴力迫害他。\par

% structure: p0009 source=h2 target=subsection reason=relative-heading-rank; base=h2

\subsection{三十八、宫廷中太监的邪恶势力}

书信的内容既是如此，恐惧与欺诈立即在全城蔓延。多少家庭遭到威胁！多少人得到重大许诺，条件是要他们对付\Person{利伯略}！多少主教目睹这些事，便藏匿起来！多少贵妇由于\Person{基督}仇敌的诽谤而退居乡间！多少苦修者成为他们阴谋的对象！多少寄居在那里并以那里为家的人，遭受他们的迫害！他们何等频繁、何等严密地看守港口和城门入口，以免任何正统教徒进城探望\Person{利伯略}！\Place{罗马}也经历了\Person{基督}仇敌的试炼，如今体会到了她先前不愿相信的事，那时她听说各城的其他教会如何遭他们蹂躏。正是那些宦官煽动这一切，攻击众人。此事中最值得注意的一点是：否认\OriginalTerm{天主之子}{Son of God}的\OriginalTerm{亚略异端}{Arian heresy}，竟从宦官那里获得支持；这些宦官，身体既无果实，灵魂又无德行，甚至无法忍受听到\Quote{子}这个名字。\Place{厄提约丕雅}的太监虽不明白自己所读的（\BibleRef{宗 8:27}），却相信了\Person{斐理伯}就\OriginalTerm{救主}{Saviour}之事教导他的话；然而，\Person{君士坦提乌斯}的宦官们却不能忍受\Person{伯多禄}的宣认，当\OriginalTerm{圣父}{Father}显示\OriginalTerm{圣子}{Son}时，他们反而转身离去，对那些说\OriginalTerm{天主之子}{Son of God}是祂真正儿子的人大发狂怒，从而声称是一个宦官的异端之说：\OriginalTerm{圣父}{Father}没有真正而真实的子。因此，法律禁止这种人进入任何教会会议；尽管如此，如今他们却把这些人视为教会事务的合格裁决者；凡他们认为好的，\Person{君士坦提乌斯}便颁下谕令，而徒有主教之名的人则与他们一道伪装。啊！谁将记述他们的历史？谁将把这些事的记录传给后代？若有人听闻，谁会真的相信呢？那些连家务都几乎不可托付的宦官（因为他们是一个喜好享乐的族类，唯一的关注便是阻挠别人拥有天性已从他们身上夺走的东西），我说这些人，如今竟在教会事务中行使权威，而\Person{君士坦提乌斯}俯从他们的意志，奸诈地密谋对付所有人，并放逐了\Person{利伯略}！\par

% structure: p0011 source=h2 target=subsection reason=relative-heading-rank; base=h2

\subsection{三十九、\Person{利伯略}对\Person{君士坦提乌斯}的讲话}

因为，皇帝多次写信给\Place{罗马}，发出威胁，派遣专员，设下计谋；随后在\Place{亚历山大}的迫害爆发后，\Person{利伯略}被拖到他面前，对他大胆直言：\Quote{停止迫害基督徒吧；不要企图借我将不虔引入教会。我们宁愿遭受一切，也不愿被称为\OriginalTerm{亚略的狂人}{Arian madmen}。我们是基督徒；不要强迫我们成为\Person{基督}的仇敌。我们也给你这劝告：不要与赐你此帝国的那一位作战，也不要以不虔取代感恩；不要迫害信靠祂的人，免得你听到这话：\BibleQuote{你用脚踢刺，是难的。}（\BibleRef{宗 9:5}）。不，我但愿你听到这些话，并能顺服，如同\Person{圣保禄}所做的那样。看，我们在这里；我们已在他们捏造指控之前来到。为此，我们急速前来，知道流放正待我们于你手中，好使我们在遭受控告之前就受苦，使所有人都能清楚看出，其他所有的人也都像我们一样受苦了；而他们所遭受的指控是仇敌的捏造，他们的一切行为不过是诽谤与虚妄。}\par

% structure: p0013 source=h2 target=subsection reason=relative-heading-rank; base=h2

\subsection{四十、\Person{利伯略}及其他人的放逐}

这是\Person{利伯略}当时所说的话，人人因此钦佩他。但皇帝并未回答，只下令将他们放逐，且将各人彼此分开，如前几次所做的一样。因为他自己在所施行的放逐中设计了这一方案，为使他的惩罚之严厉超过从前暴君和迫害者。在从前的迫害中，当时的皇帝\Person{马克西米安}曾命令将众多\OriginalTerm{宣认者}{Confessors}集体放逐，从而通过让他们彼此作伴的安慰，减轻了他们的惩罚。但这个皇帝比他更残暴；他将那些放胆直言并一同宣认的人分开，拆散那些由信仰纽带联结的人，使他们在临终时不能相见；认为身体的分离也能切断心灵的情感，又以为彼此隔绝后，他们便会忘记存于他们中间的和谐与同心。他不知道，无论各人如何与其他人分离，他仍有那与他们一同在一个身体内所宣认的\OriginalTerm{主}{Lord}同在；主也必照样预备，正如在\Person{厄里叟}先知的事例中（\BibleRef{列下 6:16}）一样，使陪伴各人的比与\Person{君士坦提乌斯}同在的士兵更多。诚然，邪恶是盲目的；因为他们以为借分隔\OriginalTerm{宣认者}{Confessors}便可折磨他们，却反而因此给自己造成极大的伤害。因为假如他们继续彼此作伴、共同居住，那些不虔之人的污秽将只在一处被宣告；但如今因将他们分离，他们便使那不虔的异端与邪恶传播开来，在各处为人所知。\par

% structure: p0015 source=h2 target=subsection reason=relative-heading-rank; base=h2

\subsection{四十一、\Person{利伯略}的失足}

凡听到他们在这场诉讼中所作所为的人，谁会不认为他们绝非基督徒，而是别的什么？当\Person{利贝里乌斯}派遣司铎\Person{欧特罗庇乌斯}和执事\Person{希拉流斯}带着书信去见皇帝时，正值\Person{路济弗尔}及其同伴宣认信仰之际，他们当场放逐了那位司铎，又剥去执事\Person{希拉流斯}的衣服，鞭打他的背，然后也将他放逐，并向他叫嚷：\Quote{你为什么不反抗\Person{利贝里乌斯}，反而替他送信？}\Person{乌尔萨奇乌斯}和\Person{瓦伦斯}，连同支持他们的阉人，是这场暴行的主谋。执事在受鞭打时赞美主，记起祂的话：\BibleQuote{我将我的背转给打击我的人}\BibleRef{依 50:6}；但他们一边鞭打他，一边嘲笑戏弄他，毫不以侮辱一个\OriginalTerm{肋未人}{Levite}为耻。的确，当他不断赞美天主时他们却发笑，这正是他们一贯的作风；因为忍受鞭打是基督徒的本分，鞭打基督徒则是\Person{比拉多}或\Person{盖法}般的暴行。这样，他们起初试图败坏罗马教会，想要将不虔敬引入其中，如同引入其他教会一样。但\Person{利贝里乌斯}被放逐两年后屈服了，因惧怕死亡的威胁而签了字。然而，这仅仅表明他们的暴行，以及\Person{利贝里乌斯}对\OriginalTerm{异端}{heresy}的憎恨和对\Person{亚大纳削}的支持，只要他被容许自由抉择的话。因为人在酷刑逼迫下，做出违背初衷的事，这不应被视为恐惧中人的自愿行为，而应归咎于施加酷刑者。然而，他们用尽一切办法支持其\OriginalTerm{异端}{heresy}，而各地教会的人民则坚守所学到的信仰，等候他们的教师归来，并谴责\OriginalTerm{敌基督的异端}{Antichristian heresy}，所有人都躲避它，如同躲避蛇蝎。\par

\SourceRecord{Translated by M. Atkinson and Archibald Robertson.；Nicene and Post-Nicene Fathers, Second Series；Vol. 4.；Edited by Philip Schaff and Henry Wace.；Buffalo, NY: Christian Literature Publishing Co.,；1892.；<http://www.newadvent.org/fathers/28155.htm>.}

% structure: p0001 source=h1 target=section reason=page-title

\section{亚略派史，第六部}

\Signature{圣亚大纳修}\par

% structure: p0003 source=h2 target=subsection reason=relative-heading-rank; base=h2

\subsection{四十二、迫害与和西的失足}

然而，尽管他们做了这一切，这些不虔敬的人认为，只要伟大的\Person{和西}逃脱了他们的邪恶阴谋，他们就一无所成。于是，他们现在又把怒火扩展到了这位伟大的老人身上。他们不以他是主教们的父亲为耻；也不顾念他曾是一位\OriginalTerm{宣认者}{Confessor}；更不尊重他主教任期之长，他任主教已超过六十年；他们却把一切置之度外，只着眼于自己异端的利益，真如那些不敬畏天主、也不尊重人的人。\BibleRef{路 18:2} 因此，他们去见\Person{君士坦提乌斯}，再次用了如下论调：\Quote{我们已经尽力了；我们流放了罗马的主教；在他之前还有大量其他的主教，并使各处充满恐慌。但如果\Person{和西}还在，您这些严厉措施对我们就算不得什么，我们的成功也决不会更稳固。只要他在自己的地方，其余的人也就留在各自的教会，因为他能以论证和信德说服众人反对我们。他是大公会议的主席，他的信函处处受人遵从。是他提出了\WorkTitle{尼西亚信经}，并到处宣布亚略派是异端者。因此，如果容许他留下，流放其他人就毫无用处，我们的异端必将覆灭。那就开始迫害他吧，不要顾惜他，虽然他年事已高。我们的异端不知道尊重老年人的白发。}\par

% structure: p0005 source=h2 target=subsection reason=relative-heading-rank; base=h2

\subsection{四十三、和西的英勇抵抗}

皇帝听到这话，不再迟延，而是了解这个人及其年高德劭，就写信召见他。这发生在他刚开始对付\Person{利伯里乌斯}的时候。他到来后，皇帝劝说他，并用那些他以为也能欺骗其他人的惯常论调催促他，要他签名反对我们，并与亚略派共融。但这位老人几乎听不下去，因皇帝竟敢提出这样的建议而感到悲伤，于是严厉地斥责了他，并在得到皇帝同意后，退回到自己的地方和教会。然而异端者仍在抱怨，并挑唆皇帝继续（他还有宦官提醒他并进一步催促他），皇帝又写了威胁性的信；但\Person{和西}虽然忍受他们的侮辱，却丝毫不动于他们对他阴谋的恐惧，他坚定于自己的目标，如同将信德之屋建立在磐石上的人，他大胆反对异端，视信中提出的威胁不过如同雨滴和风袭。尽管\Person{君士坦提乌斯}频繁写信，有时奉承地称他为父，有时威胁并列举那些被流放者的名字，还说：\Quote{难道你要继续做唯一反对异端的人吗？听从劝告，签名反对\Person{亚大纳修}吧；因为凡签名反对他的，就是与我们一同拥护亚略派的事业；} 然而\Person{和西}仍无所畏惧，且在忍受这些侮辱期间，写了一封回信，内容如下。我们读过这封信，它被置于末尾。\par

% structure: p0007 source=h2 target=subsection reason=relative-heading-rank; base=h2

\subsection{四十四、\Quote{和西致君士坦提乌斯皇帝，在主内请安。}}

我最初在您祖父\Person{马克西米安}时代兴起教难时便是一名\OriginalTerm{宣信者}{Confessor}；倘若您要迫害我，现在我也已准备忍受一切，而不愿流无辜者的血、背叛真理。但我不能赞同您以这种威胁的方式写信。请停止这样写，不要接受\Person{亚略}的主张，不要听从\Place{东方}的那些人，也不要相信\Person{乌尔撒奇}、\Person{瓦伦斯}及其同党。因为他们无论断言什么，都不是为了\Person{亚大纳削}的缘故，而是为了他们自己的异端。请相信我的话，\Person{君士坦提乌斯}啊，我年纪足以作您的祖父。我曾出席\Place{撒尔底加}会议，那时您和您那令人怀念的弟弟\Person{君士坦斯}将我们全体召集在一起；出于我本人的意愿，当\Person{亚大纳削}的敌人们来到我居住的教会时，我向他们提出挑战，若他们掌握任何对他不利的事，便可当众宣布；我愿他们放心，并期待一切事情都将获得公正的裁判。我这样做了不止一次，我请求他们，若他们不愿出席全体会议，至少可以单独前来见我；我还应许他们，若他被证明有罪，我们定会将他开除；但若查明他无可指摘，且证实他们是诽谤者，那么，如果他们随后仍拒绝与他共融，我将说服他随我一同前往\Place{西班牙}。\Person{亚大纳削}愿意接受这些条件，对我的提议未表示任何反对；但他们完全不信任自己的理由，不肯同意。又有一次，当您下诏召见他时，\Person{亚大纳削}来到了您的宫廷，那时他的敌人正在\Place{安提约基}，他便请求将他们全体或分别召来，为要使他们要么定他的罪，要么自己被定罪，并且要么当着他的面证实他们所描述的罪名，要么就停止在他缺席时控告他。对这提议，您也不肯听从，他们同样予以拒绝。那么，您为何仍听从那些说他坏话的人呢？您怎能容忍\Person{瓦伦斯}和\Person{乌尔撒奇}，尽管他们已经撤回前言，并就其诽谤作出了书面承认呢？因为他们假装是被迫承认的，这并非事实；当时并没有士兵在旁胁迫他们；您的弟弟对此事毫不知情。不，在他的治理下从未发生过如今这样的事；愿天主禁止。其实，他们是自愿上到\Place{罗马}，当着主教和众司铎的面写下了撤销声明，还事先致信\Person{亚大纳削}，言辞友善而平和。倘若他们声称曾遭受暴力，并承认这是一桩恶事，而您也对此表示不赞同，那么请您自己也停止使用暴力；不要再写信，不要再派遣\OriginalTerm{宫廷官员}{Counts}；而要释放那些已被流放的人，免得当您抱怨暴力之时，他们却变本加厉地施行更大的暴行。\Person{君士坦斯}何曾做过这样的事？有哪位主教遭受过流放？他何时曾充当教会审判的仲裁者？他的哪位\OriginalTerm{宫廷侍从}{Palatine}曾强迫人签名反对某人，以致\Person{瓦伦斯}及其同党竟能如此声称呢？我恳求您，停止这些行为，要记得您也是一个会死的人。要畏惧审判之日，并为此保持自己的纯洁。不要插手教会的事务，也不要就这些事向我们发号施令；而应向我们学习这些事务。天主将国权交在您手中；祂将祂的教会的事务托付给了我们；正如有人若想从您手中窃取帝国，便是违抗天主的命令，同样，您也当畏惧，免得您擅自接管教会的治理，而犯下大罪。经上记载：\BibleQuote{凯撒的，就应归还凯撒；天主的，就应归还天主。}\BibleRef{玛 22:21}因此，我们既不可行使属世的统治权，您，陛下，也没有焚香的权力。我出于对您救恩的关切而写了这些话。关于您信中的主题，我的决定如下：我不会与亚略派联合；我绝罚他们的异端。我也不会签名反对\Person{亚大纳削}，因为我们、\Place{罗马}教会以及整个会议都已宣判他无罪。而且您本人，在知晓这一切之后，也曾派人召见他，并准许他光荣地返回他的故乡和教会。那么，您行为上如此巨大的改变，能有什么理由呢？从前是他敌人的人，如今仍是；他们如今私下散播中伤他的言论（因为他们不敢当他的面公开说明），正是您下诏召见他之前他们所说的那些话；他们来到会议时所四处散布的，也是同样的言论。正如我先前说过的，当我要求他们站出来时，他们拿不出证据；如果他们掌握任何证据，就不会如此可耻地逃跑了。那么，是谁在这么久之后说服您忘却了自己过去的信函和声明呢？请克制自己，不要受恶人的影响，免得您为了您与他们之间的私利，却为自己招致罪责。因为您在这里满足他们的愿望，将来在审判台前，却要独自为此负责。这些人想借您之手伤害他们的敌人，希望您成为他们邪恶的工具，好藉着您的帮助，在教会中播下他们可咒异端的种子。为了取悦他人而将自己投入明显的危险，并非明智之举。所以，\Person{君士坦提乌斯}啊，我恳求您停止，请听从我的劝告。这些话由我写给您是合宜的，您也不可轻视。\par

% structure: p0009 source=h2 target=subsection reason=relative-heading-rank; base=h2

\subsection{四十五、\Person{奥西乌斯}因残酷迫害而失足}

这位如\Person{亚巴郎}般的老人\Person{霍西乌斯}，实至名归，怀有如此情怀，也写了这样的信。但皇帝并未放弃他的图谋，也未停止寻找陷害他的机会；反而继续严厉地威胁他，要么使他被迫屈服，要么在他拒绝顺从时将他放逐。正如\Place{巴比伦}的官长和总督\BibleRef{达 6:5}想找\Person{达尼尔}的把柄，却除了在他天主的法律之外，找不到任何把柄；同样，如今这些不虔敬的总督也找不出任何罪状来指控这位老人（因为这位真正的\Person{霍西乌斯}及其无可指责的生活已为人所共知），唯有一项罪名，就是憎恨他们的异端。于是，他们就开始控告他；然而，其情形并不像那些人向\Person{达理阿}控告\Person{达尼尔}那样，因为\Person{达理阿}听到指控时心中忧愁，倒是像\Person{依则贝耳}诬告\Person{纳波特}，又像\Person{犹太人}鼓动\Person{黑落德}。他们说：\Quote{他不仅不肯签名反对\Person{亚大纳削}，反而因他的缘故谴责我们；他对异端的憎恨如此之深，甚至写信给其他人，劝他们宁可受死，也不可背叛真理。\InnerQuote{因为，他说，我们亲爱的\Person{亚大纳削}也是为真理而受迫害，\Person{利伯留}（罗马主教）以及所有其余的人，都遭到了背信弃义的攻击。}}当这位不虔敬的庇护者、异端的皇帝\Person{君士坦提乌斯}，听到这些，尤其听说在\Place{西班牙地区}还有其他人持有与\Person{霍西乌斯}相同的意见，就试探他们也签名；当他无法强迫他们就范时，便召来\Person{霍西乌斯}，并没有放逐他，而是将他囚禁在\Place{锡尔米乌姆}整整一年。他无天主、不圣洁、毫无亲情，不敬畏天主，不顾自己父亲对\Person{霍西乌斯}的感情，也不尊重他极高的年龄——那时他已百岁高龄；然而，这个现代的\OriginalTerm{阿哈布}{Ahab}、我们时代的第二个\OriginalTerm{贝耳沙匝}{Belshazzar}，为了不虔敬的缘故，罔顾这一切。他对老人施加强暴，紧紧禁锢他，以致最后，老人因痛苦而屈服，勉为其难地与\Person{瓦伦斯}、\Person{乌尔萨西乌斯}及其同党共融，尽管他仍不肯签名反对\Person{亚大纳削}。但即使如此，他也未曾忘记自己的职责，因为在他临近死亡之时，如同以临终遗言一般，他见证了所遭受的暴力，绝罚了亚略异端，并严厉告诫任何人不得接受它。\par

% structure: p0011 source=h2 target=subsection reason=relative-heading-rank; base=h2

\subsection{四十六、对众多主教的任意放逐}

谁若亲眼目睹这些事，或仅仅耳闻，能不大大惊奇，并向主高声呼喊说：\BibleQuote{祢要将以色列全然灭绝吗？}\BibleRef{则 11:13}谁若了解这些行径，能不有充分理由大声呼喊说：\Quote{地上发生了一件奇妙而又可怖的事；}又说：\Quote{诸天因此而惊骇，大地更是极其恐慌。}百姓的父老和信仰的导师们被带走，而邪恶的人却被引入教会？谁看到\Person{利伯留}（罗马主教）被放逐，看到伟大的\Person{霍西乌斯}——主教之父——遭受这些苦难，或者谁看到如此多的主教被逐出\Place{西班牙}及其他地区，无论他多么迟钝，又怎能看不出那些针对\Person{亚大纳削}和其他人的指控都是虚假的，完全是诽谤？正是为此，其他的人也都忍受了这一切苦难，因为他们清楚地看到，针对这些人所设的阴谋都建立在虚假之上。因为，对\Person{利伯留}有什么指控呢？对年迈的\Person{霍西乌斯}又有什么控告呢？谁是诬告\Person{保利努斯}、\Person{路西弗}、\Person{狄奥尼修斯}和\Person{优西比乌斯}的假见证呢？或者，其余那些被放逐的主教、司铎和执事们，能被加上什么罪状呢？什么也没有；天主不许。他们身上没有任何罪名，可以用来构陷他们；他们被逐一放逐，也不是出于任何指控。这是一场不虔敬对虔敬的叛乱；这是对亚略异端的热心，是假基督来临的前奏，而\Person{君士坦提乌斯}正为此预备道路。\par

\SourceRecord{Translated by M. Atkinson and Archibald Robertson.；Nicene and Post-Nicene Fathers, Second Series；Vol. 4.；Edited by Philip Schaff and Henry Wace.；Buffalo, NY: Christian Literature Publishing Co.,；1892.；<http://www.newadvent.org/fathers/28156.htm>.}

% structure: p0001 source=h1 target=section reason=page-title

\section{亚略派历史，第七部分}

\Signature{圣亚大纳削著}\par

% structure: p0003 source=h2 target=subsection reason=relative-heading-rank; base=h2

\subsection{在\Place{亚历山大}的教难}

\Emph{之后}，他在\Place{意大利}及其他地区的教会中为所欲为之后；在流放了一些人，暴力压迫了另一些人，使各处充满恐惧之后，他终于将他的暴怒，如同某种瘟疫般的祸害，转向了\Place{亚历山大}。这是基督的仇敌巧妙策划的；因为为了使众人以为许多主教都在（声明上）签了名，并使得\Person{亚大纳削}在受迫害时连一位可以申诉的主教都没有，他们便预先行动，使各处充满恐惧，以此协助他们实施计划。然而，他们因自己的愚昧，竟不知他们展示的不是主教们的慎重选择，而是他们自己所施的强暴；并且，纵使他的弟兄离弃他，他的朋友和相识远远站立，无人同情他、安慰他，但远超这一切，有天主作他的避难所，为他就足够了。因为\Person{厄里亚}在受迫害时也是独自一人，天主是那位圣者的一切。救主在这方面也给了我们榜样，祂也曾被独自留下，暴露在仇敌的阴谋之下，为教训我们：当我们受迫害、被人离弃时，不可灰心，而应寄望于祂，不可背叛真理。因为起初真理虽看似受折磨，但日后连迫害者也将承认它。\par

% structure: p0005 source=h2 target=subsection reason=relative-heading-rank; base=h2

\subsection{对\Place{亚历山大}教会的攻击}

于是他们催促皇帝，皇帝先写了一封威胁信，派人送给公爵和士兵们。书记官\Person{狄奥格尼乌斯}和\Person{依拉略}及一些宫廷官员是送信的人；他们一到，那些可怕而残酷的暴行便施于教会，就是我上文中略述过的，也是尽人皆知的，因为民众提出的抗议书已附于本史末尾，人人均可阅读。此后，在\Person{西里亚努斯}的这些行为之后，在这些暴行既已实施、且对童贞女施以暴力之后，他竟认可这种行径和加于我们的灾祸，再次写信给亚历山大元老院和民众，煽动青年，要求他们聚集起来，要么迫害\Person{亚大纳削}，要么自视为皇帝的敌人。然而，在这些指示送达之前，在\Person{西里亚努斯}闯入教会的那时起，他就已退避了；因为他记起经上记载的话：\BibleQuote{我的百姓！你们进去，关上门，隐藏片刻，直到愤怒过去。}\BibleRef{依 26:20}有一个官阶为伯爵的\Person{赫拉克略}是这封信的送信人，也是皇帝派来的一个叫\Person{乔治}的密探的先驱——因为被皇帝派来的人不可能成为主教；天主不容！事实上他的行为及他到来之前的情况已充分证明了这点。\par

% structure: p0007 source=h2 target=subsection reason=relative-heading-rank; base=h2

\subsection{\Person{君士坦提乌斯}假意尊重亡兄的虚伪}

\Person{赫拉克略}随即公布了这封信，这使写信人蒙受极大的羞辱。因为当伟大的\Person{何西乌斯}致信\Person{君士坦提乌斯}时，他无法为自己的行为转变找到任何说得过去的借口，如今却捏造了一个对他自己及其谋士更为可耻的理由。他说：\Quote{因着我对我那位具有神圣而虔诚记忆的兄弟所怀有的感情，我暂时容忍了\Person{亚大纳削}来到你们中间。}这证明他既违背了承诺，又在兄弟死后忘恩负义。然后他宣称这位兄弟，事实上也的确，\Quote{配受神圣而虔诚的记念}；然而，对于他的命令，或用他自己的话来说，他对他所怀的\Quote{感情}，即使他只是\Quote{为了}有福的\Person{君士坦斯}的\Quote{缘故}而服从，他也理应对他的兄弟公正行事，使自己既继承帝国也继承他的情感。然而，虽然他在寻求自己正当权利时，曾以\Quote{兄弟死后，遗产归于谁？}的问题废黜了\Person{韦特拉尼奥}，却为了基督仇敌的可咒骂的异端，罔顾正义的要求，对兄弟未尽孝道。不但如此，为了这异端，他甚至不愿丝毫不违背父亲的意愿；但是，凡为取悦这些不虔之人，他便假装采纳父亲的意图，而为了折磨别人，他却不顾对父亲应有的尊重。事实上，由于\Person{欧瑟伯}及其同党的诬告，他父亲曾将那位主教暂时送往\Place{高卢}，以躲避迫害者的残暴（这事由那位已故者之兄弟、有福的\Person{君士坦丁}在他们父亲死后，通过信件表明出来），但他却不肯受\Person{欧瑟伯}及其同党的劝说，派遣他们所想要的人为主教，反而阻止他们愿望的实现，并用严厉的威胁制止了他们的企图。\par

% structure: p0009 source=h2 target=subsection reason=relative-heading-rank; base=h2

\subsection{\Person{君士坦提乌斯}如何表现对父兄的敬意}

如果因此，正如他在信中所宣称的，他渴望遵循他父亲的惯例，为什么他先差遣\Person{额我略}，现在又差遣这个\Emph{吞噬仓库者}\Person{乔治}？为什么他如此竭力地设法将这些他父亲称为\OriginalTerm{波菲里派}{Porphyrians}的\OriginalTerm{亚略派}{Arians}引入\OriginalTerm{教会}{Church}，并在庇护他们的同时驱逐其他人？尽管他的父亲曾允许\Person{亚略}觐见，但当亚略发虚誓并肚腹崩裂而死时，他便失去了父亲的怜悯；他父亲在得知真相后，便判定他为\OriginalTerm{异端者}{heretic}。此外，为什么他在假装尊重\OriginalTerm{教会法典}{Canon of the Church}的同时，却下令让自己的整个行为方式与之相悖？因为哪里有一条法典规定\OriginalTerm{主教}{Bishop}应从宫廷中任命呢？哪里有一条法典允许士兵闯入教会呢？有什么传统允许伯爵和无知的宦官在\OriginalTerm{教会事务}{Ecclesiastical matters}中行使权力，并通过他们的法令公布那些拥有主教之名者的决定呢？为了这\OriginalTerm{邪恶的异端}{unholy heresy}，他犯下了种种谎言。此前，他违反父亲的意愿，再次差遣\Person{斐拉格里乌斯}出任\OriginalTerm{总督}{Prefect}，而我们现在看到了所发生的事。他也不是\Quote{为了他的兄弟的缘故}说真话。因为在他兄弟去世后，他不是一次两次，而是三次写信给主教，反复向他承诺不会改变对他的态度，并劝他鼓起勇气，不要因任何人而惊慌，而要继续平安无事地留在他的教会里。他还通过伯爵\Person{阿斯特里乌斯}和公证人\Person{帕拉迪乌斯}，向当时任\OriginalTerm{公爵}{Duke}的\Person{费利西姆斯}和总督\Person{聂斯托里乌斯}下达命令，如果总督\Person{斐理伯}或任何其他人胆敢对\Person{亚大纳削}施以阴谋，他们就要加以阻止。\par

% structure: p0011 source=h2 target=subsection reason=relative-heading-rank; base=h2

\subsection{五十二、皇帝无权统治教会}

因此，当\Person{第欧根尼}到来，\Person{叙利亚努斯}埋伏等候我们时，他、我们和民众都要求查看皇帝的信函，因为我们想，正如经上\BibleQuote{不可在君王面前妄说谎言}所载，那么，君王一旦许下承诺，就不会撒谎，也不会改变。如果他是\Quote{为了他的兄弟的缘故而顺从}，为什么在他兄弟死后却又写了那些信呢？如果他写信是\Quote{为了对他的怀念}，为什么后来却对他如此刻薄，逼迫他，并写出那样的信，假托主教们的判决，而实际上只是为了取悦自己呢？然而他的狡诈并未逃过察悉，我们手头就有证据。因为如果真有主教们的判决，皇帝何需干预？如果只是皇帝的恐吓，那么所谓的那些主教又有何必要？从创世以来，何时听说过这样的事？教会的判决何时曾需要皇帝的认可才有效？或者更确切地说，皇帝的法令何时曾被教会承认过？历史上曾召开过许多大公会议，教会也作出过许多判决，但教父们从不寻求皇帝的同意，皇帝也从不干预教会事务。宗徒\Person{保禄}在\Person{凯撒}的家中有些他认识的人，在\BibleBook{斐理伯}{书}中他也向他们问安，但他从未带他们参与教会的裁断。然而现在我们目睹了一幕新景象，这正是\OriginalTerm{亚略异端}{Arian heresy}的发明。异端者们与皇帝\Person{君士坦提乌斯}聚在一起，好让皇帝假借主教们的权威，为所欲为地运用他的权力，从而在施行迫害时能摆脱迫害者的名号；而他们也得以在皇帝政权的支持下，任意合谋毁灭那些不与他们的不虔敬同流的人。人们可以把他们的所作所为视作一场在舞台上演出的喜剧，那所谓的众主教是演员，而\Person{君士坦提乌斯}则是执行他们命令之人，他对他们许下承诺，正如\Person{黑落德}对\Person{黑落狄雅}的女儿所做的那样，他们在他面前跳舞，通过诬告达成了放逐与杀害那些真正信仰主的人的结局。\par

% structure: p0013 source=h2 target=subsection reason=relative-heading-rank; base=h2

\subsection{五十三、\Person{君士坦提乌斯}的专制干涉}

谁没有受过他们的诽谤之害呢？这些基督的仇敌图谋毁灭过谁呢？\Person{君士坦提乌斯}何曾未按他们所提的指控放逐过谁呢？他何时拒绝过倾听他们呢？最奇怪的是，他何时允许过任何人反驳他们，而不肯轻易采纳他们的证言，不论那证言是何种性质呢？哪里有教会如今享有自由敬拜基督的特权呢？如果有教会是真正虔诚的拥护者，它便陷于危险；如果它装假，它便生活在恐惧中。就地而言，处处充满了伪善和不虔敬；到处都有虔诚者和热爱基督的人（这样的人到处都有许多，一如曾有的先知们和伟大的\Person{厄里亚}），他们隐藏自己，倘若能找到像\Person{敖巴狄亚}那样忠信的朋友，他们就退避到地上的洞穴和地穴中，或流离在旷野里。这些人在疯狂中，竟仿效\Person{依则贝耳}对\Person{纳波特}的诬告，以及犹太人对救主的诬陷，来捏造对他们的诽谤；而这位皇帝，身为异端的庇护者，且欲颠倒真理，正如\Person{阿哈布}想把葡萄园变为菜园一样，他们对他的提议正合他私意，他就照他们愿望而行。\par

% structure: p0015 source=h2 target=subsection reason=relative-heading-rank; base=h2

\subsection{五十四、\Person{君士坦提乌斯}将\Place{亚历山大}教会交给异端者}

因此，如我先前所说，他放逐了那些真正的主教，因为他们不愿为迎合他的喜好而宣认不虔敬的道理；于是他如今派遣\Person{赫拉克利乌斯伯爵}去对付\Person{亚大纳削}。伯爵公开宣布了他的决定，并传达了皇帝的命令：除非他们遵从信中指示，否则将断绝他们的粮食，推倒他们的偶像，并将许多城里的官员和百姓交予某种奴役。他如此威胁之后，竟不知羞耻地大声公开宣告：\Quote{皇帝否认\Person{亚大纳削}，并已下令将教堂交给亚略派。}众人听见都惊讶不已，彼此示意，喊道：\Quote{什么！\Person{君士坦提乌斯}成了异端者吗？}这人本应脸红，却反倒变本加厉，迫使元老院议员、异教官员和偶像庙宇的看守人同意这些条件，并承诺接受皇帝派来的任何主教。当然，\Person{君士坦提乌斯}如此行事时，正严格维护着教会的法典；他本应从教会获取信函，却转而从市场索取；本应征求民众的意见，却转而询问庙宇的看守人。他意识到自己并非派遣一位主教去治理基督徒，而是为那些同意他条件的人送去一个闯入者。\par

% structure: p0017 source=h2 target=subsection reason=relative-heading-rank; base=h2

\subsection{五十五、闯入大教堂}

外邦人因此，以为顺从便可保全偶像的安全，而某些行业的人，虽然不情愿，但因畏惧他的威胁，也签了字；恰如任命将军或其他官员一般。事实上，作为异教徒，他们除了照皇帝的心意行事，还能做什么呢？然而，民众已聚集在大教堂内（因为那天是星期的第四天），次日，\Person{赫拉克利乌斯伯爵}带着埃及太守\Person{卡塔弗罗尼乌斯}、总税务官\Person{福斯提努斯}，以及异端者\Person{比提努斯}；他们一同煽动普通暴民中那些崇拜偶像的年轻人，去攻击大教堂，用石头打民众，声称这是皇帝的命令。由于散会的时间已到，大部分人已离开了大教堂，但仍有少数妇女留下，他们便照着指使去做，于是上演了一幕悲惨的景象。那少数妇女刚祈祷完起身坐下，年轻人就突然赤身而来，拿着石头和棍棒。有些妇女被这些邪恶之徒用石头打死；他们用鞭子抽打圣洁贞女的身体，撕掉她们的面纱，使她们的头暴露在外；当她们抵抗这种羞辱时，这些懦夫就用脚踢她们。这令人恐惧，极其恐惧；但随后发生的事更糟，比任何暴行都更令人难以忍受。他们知道贞女的圣洁品格，知道她们的耳朵不习惯污秽的言语，也知道她们宁愿忍受石头和刀剑，也不愿忍受淫词秽语，于是就用这种言语攻击她们。这些都是亚略派教唆那些年轻人干的，他们还嘲笑他们所说所做的一切；而圣洁的贞女和其他敬虔的妇女，如同躲避毒蛇的咬噬一般，逃离这些言语；但基督的仇敌却协助他们的恶行，甚至可能自己也说出了同样的话；因为他们喜欢那些年轻人对她们发泄的淫秽言语。\par

% structure: p0019 source=h2 target=subsection reason=relative-heading-rank; base=h2

\subsection{五十六、大教堂遭劫掠}

此后，为彻底执行所接获的命令（这正是他们急切渴望的，也是伯爵和总税务官指示他们做的），他们夺取了座位、宝座、木制的桌子，以及教堂的窗帘和所有他们能拿到的东西，搬出去在门前大街上焚烧，并在火焰上撒了乳香。唉！谁听了这些事能不流泪，或许还会掩耳，不愿聆听这叙述，认为仅仅听闻如此滔天罪行都是有害的呢？此外，他们称颂自己的偶像，说道：\Quote{\Person{君士坦提乌斯}成了异教徒，亚略派已承认我们的习俗；}因为只要能确立他们的异端，他们甚至不惜伪装成异教信仰。他们甚至准备献祭一头为 \Place{恺撒里昂} 花园拉水的小母牛；若不是因为是雌性，他们早就献祭了；因为他们说，在他們的习俗中，献祭这样的牲畜是不合法的。\par

57. 如此，不虔敬的\OriginalTerm{亚略派}{Arians}与异教徒勾结行事，以为这些事会带给我们羞辱。但天主的公义斥责了他们的罪恶，并施行了一个伟大而显著的标记，以此清楚地向众人表明：正如他们在不虔敬的行为中胆敢攻击的不是别人，正是主，同样在这些行径中，他们又试图羞辱主。这由随后发生的奇妙事件更加明显地证明了。这些放荡的青年人中有一个跑进教会，竟敢坐在宝座上；当他坐在那里时，那个可怜人用鼻音哼唱了一首淫荡的歌曲。然后他站起来试图拉动宝座，并把它拖向自己；他不知道他正在将报应引到自己身上。因为正如古时\Place{阿左托}的居民，当他们胆敢触碰那即使看都不许看的约柜时，便立时被约柜击杀，先被痔疮痛苦地折磨；同样这个妄想拖动宝座的不幸之人，将宝座拉向自己，仿佛天主的公义派遣那木头来惩罚他，他把宝座猛撞进自己的肚腹；他非但没有搬走宝座，反而一下子撞出了自己的内脏；结果宝座夺去了他的性命，而不是他搬走宝座。正如经上关于\Person{犹达斯}所载，他肚腹崩裂，肠子都流出来；随即倒跌在地，被带走，第二天就死了（\BibleRef{宗 1:18}）。另外还有一人，拿着树枝进入教会，以外邦人的方式用手挥舞，加以嘲弄，立刻就被击打失明，顿时失去视力，不知自己身在何处；当他将要跌倒时，同伴们过来拉住他的手，扶他离开那地方，到第二天，他好不容易才回过神来，却不知道由于自己的胆大妄为，他做了什么或遭受了什么。\par

% structure: p0022 source=h2 target=subsection reason=relative-heading-rank; base=h2

\subsection{五十八、\Place{亚历山大里亚}的大迫害}

外邦人看见这些事，便心生恐惧，不敢再行任何暴行；但\OriginalTerm{亚略派}{Arians}却仍不感到羞耻，反而像犹太人见了奇迹却不信从，硬着心不肯相信，又像\Person{法郎}那样心硬；他们也把希望寄托在下，即寄托于皇帝和他的太监们。他们容许外邦人，甚或是外邦人中最放纵的人，照前述方式行事；因为他们发现，那个名义上是总税务司、却生活习惯粗俗、心地放荡的\Person{福斯蒂努}，竟然愿意在这些事上与他们勾结，煽动异教徒。不仅如此，他们自己也着手作类似的事，正如他们将他们的异端杂糅了所有其他异端，他们也要与最邪恶的人分享他们的恶行。他们借他人之手所行之事，我已在上文述说；他们亲手所犯的暴行，则超出一切邪恶的界限，且比任何刽子手都更为恶毒。哪一家他们没有捣毁过？哪一户他们未曾以搜查对手为借口而洗劫？哪一座园子他们没有践踏？哪一座坟墓他们没有挖掘，假装寻找\Person{亚大纳削}，其实只是要抢夺所遇到的一切？多少人的家被查封！他们将多少人的住所内财物分给协助他们的士兵！谁没有经历过他们的恶行？谁遇见他们却不得不躲避在闹市之中？多少人因畏惧他们而离家，在旷野过夜？多少人本欲从他们手中保全财产，反而损失大半？谁即使不谙水性，却宁愿投身大海，冒着一切危险，也不愿面对他们的威胁？许多人还更换了住所，从一条街搬到另一条街，从城里搬到郊外。许多人被课以重税，无力缴纳时，便向他人借贷，只为能逃脱他们的阴谋。\par

% structure: p0024 source=h2 target=subsection reason=relative-heading-rank; base=h2

\subsection{五十九、\Person{塞巴斯提安}的暴行}

因为他们让众人感到畏惧，以皇帝的名义、并以皇帝的不悦威胁所有人，极其傲慢地对待他们。在恶行上协助他们的有：\Person{塞巴斯提安}公爵，一个\OriginalTerm{摩尼教徒}{Manichee}，也是个放荡的年轻人；还有郡守、伯爵，以及那个伪善的总税务司。许多谴责他们不虔、宣认真的贞女，被他们从家中拖出；其他走在街上的妇女，他们加以侮辱，还叫他们的年轻人揭开她们的头巾。他们甚至纵容他们那一伙的妇女随意侮辱人；圣洁而忠信的妇女尽管退避一旁，给他们让路，他们却像\OriginalTerm{酒神狂女}{Bacchanals}和\OriginalTerm{复仇女神}{Furies}一样围住她们，若无法加害她们，便视为不幸，若当天未能对她们作些恶事，便觉愁苦。总之，他们对所有人如此残忍刻毒，以致众人称他们为刽子手、杀人犯、不法之徒、闯入者、作恶者，而不用基督徒之名，甚或其他任何名称。\par

% structure: p0026 source=h2 target=subsection reason=relative-heading-rank; base=h2

\subsection{六十、\Person{欧提基}的殉道}

此外，他们效仿\OriginalTerm{西徐亚人}{Scythians}的野蛮行径，抓住了一位曾忠心服务教会的\OriginalTerm{副执事}{Subdeacon}\Person{优提基乌斯}，用皮鞭鞭打他的背，直到他濒临死亡，然后要求将他流放到矿山去；不单是任何矿山，而是\Place{斐诺}矿场，那里连被判死刑的杀人犯也难以存活几日。他们最蛮横无理的是，连几个时辰来包扎伤口也不允许，而是下令立刻将他押送，说道：\Quote{如果这样行，众人都会害怕，从此便会归向我们。}然而，过了不久，由于鞭伤疼痛，他无法抵达矿山，死在途中。他欢欣而死，获得了殉道者的荣耀。但这些恶人仍不羞愧，反而按照圣经的话：\BibleQuote{他们毫无怜悯}\BibleRef{箴 12:10}，继续作恶，如今又行了一件属魔鬼的事。当民众为他们祈求宽免\Person{优提基乌斯}、替他说情时，他们竟逮捕了四位正直的自由公民，其中一位是替乞丐洗脚的\Person{赫尔弥亚斯}；在严厉鞭打他们之后，\OriginalTerm{司令}{Duke}将他们投入监牢。然而，比\OriginalTerm{西徐亚人}{Scythians}更残忍的\OriginalTerm{亚略派}{Arians}，见他们未因鞭打而死，便控告司令，威胁说：\Quote{我们要写信给太监们，说他不照我们想要的鞭打人。}他听了害怕，不得不再次鞭打这些人；他们受着鞭打，知道自己为何受苦，也知道是谁控告了他们，只说：\Quote{我们是为真理受鞭打，但我们绝不与异端者共融：你如今随你意鞭打我们罢；天主必因此审判你。}这些不虔敬的人想让他们在监牢中受危难而死；但天主的子民看准时机，为他们向司令求情，七天后或者更久，他们被释放了。\par

% structure: p0028 source=h2 target=subsection reason=relative-heading-rank; base=h2

\subsection{六十一、虐待穷人}

但亚略派对此感到恼怒，又策划了另一件更残忍、更不虔敬的事；在众人眼中极为残酷，却非常适合他们那敌基督的异端。主命令我们要顾念穷人；他说：\BibleQuote{变卖你所有的，施舍给穷人}又说：\BibleQuote{我饿了，你们给了我吃的；我渴了，你们给了我喝的；凡你们对我这些最小兄弟中的一个所做的，就是对我做的。}然而这些人，实为反对基督者，在此事上也竟敢违抗他的旨意。当司令将教堂交给亚略派，那些困苦无依的人和寡妇无法再在其中容身时，寡妇们便坐在圣职人员受托照料她们而指定的地方。亚略派看见弟兄们甘心服事并供养她们，就迫害这些寡妇，鞭打她们的脚，并在司令面前控告那些施舍给她们的人。此事是由一个名叫\Person{达纳米}的士兵执行的。这深得\Person{塞巴斯蒂安}的欢心，因为\OriginalTerm{摩尼教徒}{Manichæans}毫无怜悯；在他们中间，向穷人施怜悯竟被视为可憎之事。于是，出现了一种新的指控理由；亚略派首次发明了一种新式的审判。人们因自己所做的善行而受审；施怜悯的人被控告，受恩惠的人遭鞭打；他们宁愿让穷人挨饿，也不愿让愿意施怜悯的人施舍。这种心态，这些现代的犹太人——他们正是这样的人——是从古时的犹太人学来的，那些犹太人看见生来瞎眼的人复明，长期瘫痪的人痊愈，反而控告赐予这些恩惠的主，并断定那些蒙受他恩惠的人是犯法者。\par

% structure: p0030 source=h2 target=subsection reason=relative-heading-rank; base=h2

\subsection{六十二、虐待穷人}

对此，谁不感到震惊？谁不诅咒这异端及其维护者？谁看不出亚略派确实比野兽更残忍？因为他们从这些不义之行中并无利可图，原非为此而行；反倒使众人都更加憎恶他们。他们想借着诡计和恐吓胁迫一些人加入他们的异端，好让他们前来共融；但结果适得其反。受害者忍受他们所加的一切，如同殉道，既不背弃也不否认在基督内的真信仰。那些教外的人目睹他们的行为，甚至连\OriginalTerm{外教人}{heathen}最后见了这些事，也诅咒他们为敌基督者、残忍的刽子手；因为人性本倾向于怜悯穷人和同情他们。但这些人连普通的人情也丧失了；他们自己若是受苦，本会希望从别人那里得到的仁慈，却不肯容许别人领受，反而利用长官，尤其是司令的严酷和权威来对付他们。\par

% structure: p0032 source=h2 target=subsection reason=relative-heading-rank; base=h2

\subsection{六十三、虐待长老和执事}

他们对长老和执事所行的，他们如何在司令和长官们的判决下，将他们放逐，又命士兵把他们的亲属从家中拖出，以及警务长\Person{戈尔哥尼}如何鞭打他们；还有（最残忍的举动）他们如何蛮横地抢夺这些人和那些已故者的食粮；这些事非言语所能形容，因为他们的残酷超越了语言的一切能力。有什么词语能配得上这样的行径？该先提及哪一件，才能让后续记载的事显得不那么骇人，而下一件又更加可怖？他们一切的图谋和不义都充满凶杀和不虔；而且他们如此不择手段、诡计多端，试图用保护他们的承诺和金钱收买来欺骗人，这样，既然不能以正当手段博得人心，就借此炫耀，以欺骗单纯的人。\par

\SourceRecord{Translated by M. Atkinson and Archibald Robertson.；Nicene and Post-Nicene Fathers, Second Series；Vol. 4.；Edited by Philip Schaff and Henry Wace.；Buffalo, NY: Christian Literature Publishing Co.,；1892.；<http://www.newadvent.org/fathers/28157.htm>.}

% structure: p0001 source=h1 target=section reason=page-title

\section{亚略派历史，第八部分}

\Signature{圣亚大纳削}\par

% structure: p0003 source=h2 target=subsection reason=relative-heading-rank; base=h2

\subsection{六十四、埃及的迫害}

有谁还会称呼他们为 \OriginalTerm{外邦人}{Gentiles} 呢？更别提称呼他们为基督徒了！有谁会认为他们的习性和情感是人性，而非野兽的呢？看看他们那残忍野蛮的行为吧！他们比公开的刽子手更无价值；比所有其他异端者更胆大妄为。他们远不如外邦人，而且与外邦人格格不入、判然有别。我从教父们那里听说过，并且我相信他们的记载是真实的：许久以前，在 \Person{马克西米安}——\Person{君士坦提乌斯}的祖父——的时代，迫害兴起时，外邦人曾藏匿我们被搜捕的基督徒弟兄，常常为此损失自己的财产，遭受监禁的试炼，为的就是不交出逃亡者。他们保护那些投奔他们避难的人，如同保护自己一样，并决心为他们冒一切风险。然而现今这些可钦佩的人——新异端的发明者——所为却完全相反；除了其背叛行径外，别无可称。他们自命为刽子手，图谋出卖所有人，并设计陷害那些藏匿他人者，将藏匿者和被藏匿者同视为敌人。他们是如此嗜杀，在作恶上如此效法 \Person{犹达斯}的邪恶。\par

% structure: p0005 source=h2 target=subsection reason=relative-heading-rank; base=h2

\subsection{六十五、\Place{巴尔卡}的\Person{塞孔杜斯}的殉道}

这些人所犯的罪行无法充分描述。我只愿说，当我执笔并且想悉数列举他们一切恶行时，心中浮现这样的念头：这个异端岂不就是\BibleBook{}{箴言}中所说水蛭的第四个女儿（见：\BibleRef{箴 30:15}）？因为在如此多的不义之行、如此多的谋杀之后，它还未曾说：\Quote{够了。}不！它还在疯狂肆虐，四处寻觅那些尚未被发现的人，而对于那些它已经伤害过的人，它还想再次加害。在夜袭之后，在随之而来的一切灾祸之后，在赫拉克里乌斯所挑起的迫害之后，他们仍不停地在皇帝面前诬告我们（且他们深信，像他们这样不虔敬的人，必会得到垂听），希望对我们施加比放逐更重的刑罚，并愿从此以后，凡不顺从他们亵渎之事的人，尽遭毁灭。于是，他们现在气焰极其嚣张。那个 \Place{五城}的 \Person{塞孔杜斯}，一个最放荡的人，连同他的同党 \Person{斯德望}，深知他们的异端可以辩护他们所犯的任何不义，便在 \Place{巴尔卡}发现了一位不愿顺应他们欲望的司铎（他也叫塞孔杜斯，与异端者同名，但信仰不同），于是他们将他踢至死。当他受此苦时，他效法那位圣者，说：\Quote{愿无人为我向人间审判官伸冤；我自有主为我伸冤，我正是为祂的原故，在他们手下忍受这些苦。}然而，他们听了这话并不动心，亦未对圣期有丝毫敬畏；因为那时正值 \OriginalTerm{四旬期}{Lent}内，他们竟如此将人踢死。\par

% structure: p0007 source=h2 target=subsection reason=relative-heading-rank; base=h2

\subsection{六十六、迫害——\OriginalTerm{亚略主义}{Arianism}的武器}

哦，新异端啊，你在不敬和恶行上集中了魔鬼的全部作为！因为事实上，它不过是一个新近发明的邪恶；虽然以前似乎也有人采纳过它的教义，但他们秘而不宣，无人知晓。然而像从洞中爬出的蛇一般，\Person{欧瑟比乌斯}和 \Person{亚略}吐出了这不敬的毒液；亚略公然亵渎神明，欧瑟比乌斯则为其亵渎辩护。然而，他却无力维持这异端，直到如前所述，他在皇帝那里找到了一位保护者。我们的教父们召开了一次 \OriginalTerm{大公会议}{Ecumenical Council}，约有三百位教父聚集一处，共同谴责了亚略异端，并一致宣明它远离且外乎教会的信仰。其支持者见自己蒙羞，且再无良善立论的依据，便想出另一种方法，试图借由外在的权力捍卫这异端。在此，人尤当惊叹其策略之新奇与邪恶，以及它如何超越所有其他异端！因为其他异端是以花言巧语之论证欺骗愚拙的人；正如宗徒所说，希腊人以美妙的言辞与说服力进攻，以看似有理的谬论制胜；犹太人则离弃神圣的圣经，正如宗徒再次所说的那样，争辩的是 \BibleQuote{无稽的传说和无穷的祖谱}（见：\BibleRef{弟前 1:4}）；而 \OriginalTerm{摩尼教徒}{Manichees} 和 \OriginalTerm{华伦提努斯派}{Valentinians} 等，则败坏圣经，按自己的虚构发表传说。但 \OriginalTerm{亚略派}{Arians} 比他们所有人都大胆，并显明了其他异端只不过是他们的妹妹，正如我所说，他们在不敬上超越了他们，效法他们所有人，尤其在邪恶上特别效法犹太人。因为正如犹太人无法证明他们指控保禄的捏造罪名时，便立即将他扭送千夫长和总督那里；同样地，这些在计谋上超越犹太人的家伙，只使用审判官的能力；若有人稍微反对他们，便被拖到总督或将军面前。\par

% structure: p0009 source=h2 target=subsection reason=relative-heading-rank; base=h2

\subsection{六十七、\OriginalTerm{亚略主义}{Arianism}因迫害而恶于其它异端}

其他异端也是如此，当真道本身以最清晰的证据驳斥它们时，它们通常保持沉默，仅仅因被定罪而不知所措。但这新近且该受诅咒的异端，当它被论证击败，被真道本身打倒并蒙羞之时，便立刻企图用暴力、鞭打和监禁来胁迫那些它无法用论证说服的人，从而承认自己远非虔诚。因为真正虔诚的本分不是强迫，而是说服，正如我先前所说。因此，我们的主亲自对众人说：\BibleQuote{若有人要跟随我}\BibleRef{玛 16:24}；并对门徒们说：\BibleQuote{难道你们也要走吗}\BibleRef{若 6:67}？祂并没有使用武力，而是让他们自由选择。然而，这个异端却与虔敬完全格格不入；因此，它除了行事与我们的救主相反之外，还能怎样呢？更何况它已招募了基督的敌人\Person{君士坦提乌斯}，这个如同假基督本身的人，作为它不信虔敬的首领。他为了这个异端，极力效仿\Person{撒乌耳}的野蛮残忍。因为当司祭们给\Person{达味}食物时，\Person{撒乌耳}下令，他们便被全部杀死，人数是三百零五；而这个人，如今所有人都回避异端，并宣认对主纯正的信仰，他却废除了一个足有三百位主教的大公会议，流放了主教们自己，并阻止人民实践虔敬，阻止他们向天主祈祷，禁止他们的公共集会。正如\Person{撒乌耳}推翻了司祭城\Place{诺布}，这个人，在邪恶上更进一步，已将教会交给了不虔者。正如他在真司祭面前抬举控告者\Person{多厄格}，并迫害\Person{达味}，听信齐弗人；这个人也偏爱异端者超过虔诚者，并仍然迫害那些躲避他的人，听信他自己的宦官们，他们诬告正统派。他没有意识到，他为亚略派异端所做或所写的一切，都涉及对救主的攻击。\par

% structure: p0011 source=h2 target=subsection reason=relative-heading-rank; base=h2

\subsection{六十八、\Person{君士坦提乌斯}比\Person{撒乌耳}、\Person{阿哈布}和\Person{比拉多}更恶劣。他过去对自己亲属的所作所为}

\Person{阿哈布}本人对待天主的司祭，都不像此人对待主教们那样残忍。因为至少在\Person{纳波特}被谋杀后，他良心受到刺痛，见到\Person{厄里亚}时\BibleRef{列上 21:20}感到害怕；但此人既未尊敬伟大的\Person{何西乌}，在流放如此多主教后，也不感到厌倦或良心刺痛；反而像另一个法郎，越受打击，心越刚硬，日复一日地想象更大的邪恶。他极不义的事例，最为突出的是以下一桩。恰巧当主教们被判处流放时，其他一些人也因谋杀、叛乱或偷窃的罪名受到宣判，各按其罪。几个月后，他受人请求，便释放了这些人，如同\Person{比拉多}释放\Person{巴辣巴}一样；但对基督的仆人，他不仅拒绝释放，甚至判他们在流放地受更无情的惩罚，向他们证明自己是\Quote{不死的邪恶}。对其他人，他因性情相投而成为朋友；但对正统派，他却因他们真正信仰基督而成为仇敌。从此事岂非很明显：古时犹太人要求释放\Person{巴辣巴}，并将主钉死在十字架上，所行无非这些现今与\Person{君士坦提乌斯}一道的基督仇敌正在做的事？不，他甚至比\Person{比拉多}更苦毒。因为\Person{比拉多}察觉\BibleRef{玛 27:24}此事不义，就洗手；而此人却在流放圣徒时，对他们越来越咬牙切齿。\par

69. 但是，当他被引入邪恶的错误后，对主教们如此残忍，这又何足为奇呢？因为普通的人情都不能促使他放过自己的亲属。他杀了他的叔伯；清除了他的堂表亲；他不同情他岳父的遭遇，即使他娶了他的女儿，也不同情他的亲戚；他总对所有人违背誓言。同样，他以不敬虔的方式对待他的兄弟；如今他装作要修建他的坟墓，却将其兄弟一直保护至死、并抚养长大的未婚妻\Person{奥林匹亚丝}，交给了蛮族。更有甚者，他企图撇开他的意愿，尽管他自诩为他的继承人；因为他所写的信，措辞是任何稍有常识的人都会感到羞耻的。但当我比较他的信件时，我发现他并无常识，他的心思完全受他人意见的支配，根本毫无己见。\Person{撒罗满}说：\BibleQuote{如掌权者听信谎言，他的一切仆人都是恶人}\BibleRef{箴 29:12}。此人用行为证明自己正是这样一个不义者，而他周围的人也都是恶人。\par

% structure: p0014 source=h2 target=subsection reason=relative-heading-rank; base=h2

\subsection{七十、\Person{君士坦提乌斯}的反复无常}

那么，他既是这样的人，又喜欢这样的同伴，他怎能设计出任何公正或合理之事呢？他被他那些追随者的恶行所纠缠，这些人确实迷惑了他，或者更确切地说，将他的理智践踏在脚下。因此，他时而下令写信，随即又后悔写了它们，后悔之后又怒火中烧，然后再次哀叹自己的命运，并且犹豫不决，显露出一颗毫无理智的灵魂。他既是这样的品格，人们本当对他表示怜悯，因为在自由的外表和名义下，他成了那些拖着他去满足自己邪恶欲望之人的奴隶。总之，正如圣经所言，他因自己的愚昧和反复无常，愿意迎合他人的心意，就自行投入了被定罪、将来在审判中被火烧尽的结局；他既同意做他们想要的一切，又在他们谋害主教、对教会施加权威的图谋上满足他们。看哪，如今他又扰乱了所有\Place{亚历山卓}、\Place{埃及}和\Place{利比亚}的教会，并公开下令，要将天主教会及信仰的主教们赶出他们的教堂，将他们全部交给\OriginalTerm{亚略}{Arian}学说的信奉者。那位将军开始执行这项命令；主教们立刻被锁链绑缚遣送，长老和隐修士们在几乎被鞭打得死去活来之后，也被铁链捆绑。各地一片混乱；整个\Place{埃及}和\Place{利比亚}都处于危险之中，民众对这不公的命令感到愤慨，并从中看到了\OriginalTerm{敌基督}{Antichrist}来临的准备，眼见自己的财产被他人掠夺并交在异端者手中。\par

% structure: p0016 source=h2 target=subsection reason=relative-heading-rank; base=h2

\subsection{七十一、这邪恶前所未有}

何时曾听说过这样的邪恶？即使在教难时期，何时曾发生过如此恶劣的罪行？昔日迫害者都是外教人，但他们没有把偶像带入教堂。\Person{芝诺比娅}虽是个犹太女人，且支持\Person{撒摩撒塔的保禄}，但她也没有把教堂交给犹太人用作会堂。这是一桩新的恶行。这不单是教难，而是远超教难，它是敌基督来临的前奏和准备。即便承认他们捏造了虚假的罪名控告\Person{亚大纳削}和其他被放逐的主教，但与他们后来的行径相比，这又算得了什么？他们能对整个\Place{埃及}、\Place{利比亚}和\Place{五城地区}提出什么指控？因为他们已不再针对个人施展阴谋，那样的话他们或许还能编造谎言；而是整体攻击所有人，以致如果他们仅仅想对他们捏造指控，他们就必然被定罪。如此，他们的邪恶蒙蔽了他们的理智（\BibleRef{智 2:21}）；他们毫无理由地要求驱逐整个主教团，借此表明他们对\Person{亚大纳削}和其他被放逐的主教所捏造的罪名是虚假的，而且是出于唯一的目的，即支持那可诅咒的、与基督为敌的亚略派异端。此事如今已不再隐藏，而是对所有人来说都极其明显了。他下令将\Person{亚大纳削}驱逐出城，并把教堂交给了他们。与\Person{亚大纳削}在一起的、由\Person{伯多禄}和\Person{亚历山大}所任命的长老和执事，也被驱逐、流放；而那些真正的亚略派，不是因环境引起的任何嫌疑，而是因异端本身最初就由主教\Person{亚历山大}与\Person{亚略}本人一同被驱逐的——\Place{利比亚}的\Person{塞昆都}、\Place{亚历山卓}的\Person{客纳罕人欧左伊乌}、\Person{犹利乌}、\Person{阿孟}、\Person{马尔谷}、\Person{依肋内}、\Person{佐西莫}、\Person{又名佩利孔的撒拉匹翁}，以及\Place{利比亚}的\Person{西辛尼乌}和他手下的年轻人，都是与他共同行不虔之事的伙伴——这些人占据了教堂。\par

% structure: p0018 source=h2 target=subsection reason=relative-heading-rank; base=h2

\subsection{七十二、埃及主教的放逐}

将军\Person{塞巴斯蒂安}写信给各地的总督和军事长官；真正的主教遭受迫害，而那些信奉不虔敬教义的人被扶植取代他们。他们将那些在圣职上已年迈、并蒙主教\Person{亚历山大}祝圣多年的主教流放：\Person{阿摩尼乌斯}、\Person{赫尔墨斯}、\Person{阿纳甘福斯}和\Person{马尔谷}被送到\Place{上奥阿西斯}；\Person{穆伊斯}、\Person{普塞诺西里斯}、\Person{尼拉蒙}、\Person{普莱内斯}、\Person{马尔谷}和\Person{阿特诺多鲁斯}被送到\Place{阿摩尼亚卡}，其唯一目的就是让他们在穿越沙漠的途中丧命。他们虽身患疾病，却未获丝毫怜悯；事实上，由于身体虚弱，旅途极为艰难，不得不躺在担架上被人抬着，他们的病势如此危重，以致连殡葬用品都随身携带。其中一人确实死了，但他们甚至不准将尸体交还其亲友安葬。怀着同样的用心，他们还把主教\Person{德拉孔提乌斯}流放到\Place{克利斯马}附近的荒漠，把\Person{斐洛}流放到\Place{巴比伦}，把\Person{阿德尔斐乌斯}流放到\Place{底比斯}的\Place{普西纳布拉}，把长老\Person{希厄拉克斯}和\Person{狄奥斯科鲁斯}流放到\Place{塞伊尼}。他们也将\Person{阿摩尼乌斯}、\Person{阿加托斯}、\Person{阿加托德蒙}、\Person{阿波罗尼乌斯}、\Person{欧洛吉乌斯}、\Person{阿颇罗}、\Person{帕弗努提乌斯}、\Person{加犹}和\Person{弗拉维乌斯}这些古老的主教，以及主教\Person{狄奥斯科鲁斯}、\Person{阿摩尼乌斯}、\Person{赫拉克利德斯}和\Person{普萨伊斯}一同驱逐；其中一些人被抓去采石场做苦工，另一些人则遭受迫害，意图杀害，还有许多人的财物被劫掠。他们还流放了四十名平信徒，以及若干先前曾被迫置身火中的贞女；他们用棕榈树枝狠狠鞭打她们，以致其中一些人在熬过五天后死去，另一些人则因刺入肢体的荆棘而不得不接受外科治疗，所遭受的痛苦比死亡还难受。然而，最令所有理智健全之人感到恐怖、却是这些恶徒典型行径的是：当那些贞女在鞭挞中呼求\OriginalTerm{基督}{Christ}之名时，他们竟越发狂怒地向她们咬牙切齿。更有甚者，他们不肯将死者的尸体交给亲友安葬，反而藏匿起来，似乎要使谋杀不为人知。然而纸包不住火，全城都察觉到了，众人视他们为刽子手、罪犯和强盗，纷纷远离他们。此外，他们还捣毁隐修院，试图将隐修士投入火中；他们劫掠房屋，闯入一些自由公民的家宅，那里有主教存放的财宝，他们将其抢掠一空。他们鞭打寡妇的脚掌，阻止她们领取施舍。\par

% structure: p0020 source=h2 target=subsection reason=relative-heading-rank; base=h2

\subsection{七十三、亚略派提名者的特征}

亚略派所行的邪恶就是如此；至于他们更深的不虔敬行为，谁能听了不战栗呢？他们使这些德高望重的老者和年迈的主教遭受流放，如今却任命放荡的异教青年取代他们，这些人甚至还不是\OriginalTerm{慕道者}{Catechumens}，却被他们想一举提升至最高的尊位。另有些人因被指控犯有重婚罪，甚至更恶劣的罪行，却因他们拥有的财富和世俗权力而被任命为主教，仿佛从市场上被遣送出去，只因为他们给了黄金。现在，民众遭遇了更为可怕的灾难。因为当他们拒绝这些依附亚略派的雇佣之辈——这些人与他们如此格格不入——他们便遭到鞭打，被剥夺公民权，被将军（他本人是\OriginalTerm{摩尼教徒}{Manichee}，做这一切毫不手软）关进监牢，目的在于迫使他们不再寻找自己的主教，而不得不接受他们所憎恶的人，这些人如今又犯下了他们从前在自己的偶像中所行的那类丑剧。\par

% structure: p0022 source=h2 target=subsection reason=relative-heading-rank; base=h2

\subsection{七十四、\Person{君士坦提乌斯}的主教任命——反基督者的标记}

每个正直的人目睹或听闻这些事，见到这些不虔敬者的傲慢和极端不义，怎能不痛哭呢？\Quote{义人代替不敬虔者叹息。}经历了所有这一切，如今这种不虔敬已猖狂到极点，谁还敢称这个\Person{科斯蒂留斯}为基督徒，而不称之为反基督者的影像呢？还有什么反基督者的标记尚付阙如？他怎能不被视为那一位？或后者怎能不被认为就是像他这样的人？亚略派和外邦人不是在\Place{凯撒勒奥}的大教堂中，仿佛奉\OriginalTerm{基督}{Christ}之命，献上那些祭品，并且口出亵渎基督之言吗？达尼尔的神视不是这样描述反基督者的吗？\BibleQuote{他要向至高者说亵渎的话，折磨至高者的圣徒，想改变节期和律法}\BibleRef{达 7:25} 现在除了\Person{君士坦提乌斯}，还有谁曾企图做这些事呢？他确实就是反基督者将要成为的那种人。他支持这不虔敬的异端，从而向至高者说出亵渎的言语；他流放主教，从而向圣徒发动战争；尽管他行使这权力不过片时，且为自取灭亡。此外，他在邪恶上超越了前人，发明了一种新的迫害方式；在他推翻三位国王，即\Person{维特拉尼奥}、\Person{玛格嫩提乌斯}和\Person{加卢斯}之后，他立即着手庇护不虔敬之事；他骄妄无度，竟胆敢如巨人般将自己抬高，反对至高者。他妄想改变律法，通过违背主借其宗徒传予我们的规定，变更教会的习俗，并发明一种新的任命方式。因为他从远隔五十天路程的异地，派遣由士兵护送的主教，前往不愿接受他们的民众那里；他们带来的不是与当地民众相识的介绍，而是给地方长官的恐吓信函和命令。就这样，他从\Place{卡帕多细亚}派\Person{格列高利}去了\Place{亚历山大里亚}；他把\Person{革尔米尼乌斯}从\Place{基齐库斯}调往\Place{西尔米乌姆}；他把\Person{凯克洛庇乌斯}从\Place{劳狄刻亚}调往\Place{尼科美底亚}。\par

% structure: p0024 source=h2 target=subsection reason=relative-heading-rank; base=h2

\subsection{七十五、\Person{乔治}抵达\Place{亚历山大里亚}，以及\Person{君士坦提乌斯}在\Place{意大利}的行动}

他再次从\Place{卡帕多细亚}调遣一个名叫\Person{奥森提乌斯}的人前往\Place{米兰}，此人与其说是基督徒，不如说是闯入者；在将当地主教\Person{狄奥尼修斯}——一位虔诚事主的人——因其对基督的热心而流放之后，他命此人留在那里。但此人甚至不通晓拉丁语，除不虔之外一无所长。如今，一个名叫\Person{乔治}的卡帕多细亚人，此人曾任\Place{君士坦丁堡}的仓储承包商，因侵吞了所收的一切款项而被迫逃亡，皇帝命他率领军容威严进入\Place{亚历山大里亚}，并以将军的权威为后盾。接着，他发现一个初学者\Person{埃皮克提图}，一个胆大妄为的年轻人，因看出他预备行恶，便喜爱他；借着此人，他对那些他想陷害的主教们施展计谋。因为他预备好做皇帝所愿的一切；皇帝于是利用他的帮助，在\Place{罗马}行了一件奇怪的事，一件确实与\OriginalTerm{反基督}{Antichrist}的恶毒相似的事。在皇宫内而非教会中做准备，并差派大约三个他自己的太监来代替民众，然后他强迫三个卑鄙的密探（因为不能称他们为主教），在皇宫内祝圣了一个名叫\Person{菲利克斯}的人为主教，那人配得上他们。因为民众看出异端者的恶行，就不容他们进入教会，并远远避开他们。\par

% structure: p0026 source=h2 target=subsection reason=relative-heading-rank; base=h2

\subsection{七十六、\Person{君士坦提乌斯}对主教们的暴虐流放}

如今，还缺什么使他成为\OriginalTerm{反基督}{Antichrist}呢？或者\OriginalTerm{反基督}{Antichrist}来临时所能行的，还有什么比这人所行的更多呢？\OriginalTerm{反基督}{Antichrist}来临时，岂不会发现道路已被此人铺平，使他能轻而易举地欺骗百姓吗？再者，他将裁决案件的权利归于自己，并把这些案件提交法庭而非教会，且亲自主持。说起来奇怪，当他看到控告者无言以对时，他便自己担当控告者，如此使受害方因他所施的暴力而无法自辩。在针对\Person{亚大纳削}的诉讼中他正是这样做的。因为当他看到主教们\Person{保利努斯}、\Person{路西弗}、\Person{尤西比乌斯}和\Person{狄奥尼修斯}的勇气，以及他们如何基于\Person{乌尔撒西乌斯}和\Person{瓦伦斯}的悔改声明，驳斥那些反对主教的人，并建议不应再相信\Person{瓦伦斯}及其同伙，因为他们已收回了现在所断言的话，这时他立即站起来说：\Quote{我现在是\Person{亚大纳削}的控告者；你们必须因我而相信这些人所说的话。}接着，他们便说：\Quote{然而，当被告不在场时，你怎能成为控告者呢？如果你是控告者，而他却不在此，因此无法受审。况且此案并非\Place{罗马}之事，好让你作为皇帝而受信任；这是关系主教的事；因为审判应当在控告者和被告之间平等地进行。此外，你怎能控告他呢？因为你不可能亲临见证一位距你如此之远的人的行为；而且如果你所说的只是你从这些人听来的，你也应当相信他所说的话；但如果你不信他，却信他们，显然他们是为了你的缘故才这样说，控告\Person{亚大纳削}只是为了取悦于你？}当他听到这番话，认为他们这些完全真实的话是对自己的侮辱，便将他们流放；他因对\Person{亚大纳削}发怒，便以更为凶暴的措词写信，要求他应遭受如今已临到他的待遇，要求将各教会交给\OriginalTerm{亚略派}{Arians}，并允许他们为所欲为。\par

% structure: p0028 source=h2 target=subsection reason=relative-heading-rank; base=h2

\subsection{七十七、\Person{君士坦提乌斯}——\OriginalTerm{反基督}{Antichrist}的先驱}

这些行为确实可怖，且比可怖更甚；然而对于那扮演\OriginalTerm{反基督}{Antichrist}角色的人，却是相称的举止。谁看见他率领那些他声称为主教的人，主持教会案件，不会理所当然地惊呼：这岂不是\Person{达尼尔}所说的\BibleQuote{那荒凉的可憎之物}吗？\BibleRef{达 9:27} 因为他披上基督信仰的外衣，进入圣地，站立其中，却摧毁各教会，违犯教会法典，并强制推行他自己的法令。如今还有谁敢说这是基督徒的平安时期，而非迫害时期呢？这确实是一场迫害，一场前所未有的迫害，除非那\BibleQuote{不法之子}\BibleRef{得后 2:8}再来，恐怕无人能再掀起；基督的这些仇敌们正在显明他，因为他们已在自己身上呈现了\OriginalTerm{反基督}{Antichrist}的形象。因此，我们尤当清醒，免得这异端已达到如此厚颜无耻的高峰，并如\BibleBook{}{箴言}所记，如同\BibleQuote{毒蛇的毒液}\BibleRef{箴 23:32}般扩散开来，且宣讲与救主相反的教义；免得，我是说，这就是那\BibleQuote{背离}\BibleRef{得后 2:3}，此后他必要显现，而\Person{君士坦提乌斯}确实是他的先驱。否则，他为何对虔敬者如此疯狂？为何将其视为自己的异端而争辩，称凡不遵从\Person{亚略}的狂妄、不欣然接受基督仇敌的指控、不羞辱如此众多可敬大公会议的人为他的仇敌？他为何命令将各教会交给\OriginalTerm{亚略派}{Arians}？岂不是为此：当那另一位来临时，他便可由此找到进入教会之路，并收纳那为他预备这些地方的人为己所有吗？因为由\Person{亚历山大}及其前任\Person{阿基拉斯}，并在他之前的\Person{伯多禄}所祝圣的古老主教们已被逐出；而由士兵同伴所提名的人被引入；他们只提名那些承诺采纳其教义的人。\par

% structure: p0030 source=h2 target=subsection reason=relative-heading-rank; base=h2

\subsection{七十八、\OriginalTerm{梅勒提派}{Meletians}与\OriginalTerm{亚略派}{Arians}的联盟}

对于麦勒提安派来说，这是一个容易遵从的提议；因为他们中的大部分，甚至全部，从未受过宗教教育，也不知道在基督内\Quote{健全的信仰}是什么，完全不懂什么是基督教，也不了解我们基督徒拥有哪些经书。因为他们中有些人来自偶像崇拜，有些人来自元老院或首要的民政职务，为了可悲的免役和得到的庇护，就贿赂了在他们之先的麦勒提安派人士，甚至在还未受教之前就被提升到如此尊位。即便他们假称曾受过指教，但在麦勒提安派当中又能获得怎样的指教呢？事实上，他们甚至根本无意假装受教，径直前来，立刻被称为主教，就像孩子得到一个名字那样。这样的人，自然不会认为这事有多大关系，甚至不会想到虔敬与不虔敬有所不同。因此，他们轻易而迅速地由麦勒提安派变成了亚略派；如果皇帝命令他们改奉其他任何信仰，他们也会随时再变。他们对真虔敬的无知，使他们很快就顺从了那占上风的愚昧，以及最先传授给他们的东西。因为只要他们能免役，能得到人的庇护，随着各种风浪和暴风摇摆不定对他们来说根本不算什么；他们恐怕也不会顾忌再变回以前的样子，甚至变回他们作异教徒时的模样。无论如何，这些性情如此易变的人，将教会视为世俗的元老院，并且像异教徒一样怀着偶像崇拜的心态，却披戴上了救主的尊名，在这个名号下他们玷污了整个埃及，使亚略异端的名声在那里广为人知。因为埃及从前乃是唯一勇敢地维护正统信仰宣认的地方；所以这些不信真道的人努力在那里也要引入猜忌，或者说，不是他们，而是那煽动他们的魔鬼，为的是当他的先驱反基督者来临时，能发现埃及的教会也成了他自己的，麦勒提安派已经接受了他那套原则的教导，并认出自己在他们里面已经成形了。\par

% structure: p0032 source=h2 target=subsection reason=relative-heading-rank; base=h2

\subsection{七十九、麦勒提安派的行为与亚历山大的基督徒的对比}

这就是君士坦提乌斯所发出的不义命令造成的结果。民众方面表现出欣然准备殉道的热忱，以及对这至极不虔异端的愈发憎恨；然而为他们的教会而悲伤的哀叹和呻吟从各处爆发出来，他们向上主呼号说：\BibleQuote{主啊，求祢怜恤祢的子民，不要将祢的产业交给敌人，任其辱骂}\BibleRef{岳 2:17}；但求速速拯救我们脱离不法之人的手。因为看啊，\Quote{他们毫不顾惜祢的仆人，反而在为反基督者预备道路}。麦勒提安派绝不会抵抗他，也不会关心真理，更不会以否认基督为恶事。这些人从未真诚地前来接近圣言；就像变色龙一样，呈现出各种不同的外貌；他们是任何愿意利用他们的人的雇佣者。他们不以真理为目标，反倒喜爱眼前的享乐胜过真理；他们只知说：\BibleQuote{我们吃吧，喝吧，因为明天我们就要死了}\BibleRef{格前 15:32}。这样的宣认和无信的态度，更像是伊壁鸠鲁派的伶人，而不是麦勒提安派。但是，我们救主的忠仆、那些真心相信的真主教，他们活着不是为了自己，而是为了主；他们忠信地信仰我们的主耶稣基督，并且，如我以前说过的，知道那些被拿来攻击真理的指控是虚假的，显然是出于拥护亚略异端的目的而捏造的（因为乌尔撒奇和瓦伦斯的翻供，使他们识破了针对亚大纳削而编造的诽谤，这些诽谤的目的在于排除他，并将那些与基督为敌者的不虔之事引进各教会）；这些人，我说，既然察觉了这一切，身为真理的捍卫者和宣扬者，宁愿选择忍受侮辱，被放逐，也不愿签署反对他的文件，不与亚略的狂人相通。他们没有忘记自己曾教导别人的功课；是的，他们清楚地知道，背叛者等待的是极大的羞辱，而宣认真理者等待的却是天国；轻率的人和畏惧君士坦提乌斯的人得不到任何好处；但那些在此世忍受磨难的人，就像水手在风暴过后抵达平静的港湾，像角斗士在搏斗之后领受桂冠一样，这些人在天上将获得巨大而永恒的喜乐与欢欣——就像若瑟在经历了那些磨难之后所获得的；就像伟大的达尼尔在他受试探和朝臣们多方阴谋陷害之后所获得的；就像保禄现在所享有的，由救主亲手加冕；就像天主的子民各处所期望的。他们看到了这些事，就不会意志薄弱，反倒信德更为坚固，热忱日益增加。他们深知异端者的诽谤和不虔，因此定断迫害者的罪，并且在内心中与那些受迫害的人同跑一条路，好使他们也能获得宣认的荣冠。\par

% structure: p0034 source=h2 target=subsection reason=relative-heading-rank; base=h2

\subsection{八十、与异端者分离的义务}

针对这可憎的\OriginalTerm{假基督的异端}{antichristian heresy}，我们本可以说得更多，并用许多论证证明\Person{君士坦提乌斯}的行径正是\OriginalTerm{假基督}{Antichrist}来临的前兆。然而，正如先知所说（\BibleRef{依 1:6}）：\BibleQuote{从脚掌直到头顶，没有一处健全}，它充满了一切污秽和一切不虔，以致其名号当如狗呕吐之物或蛇毒一般被躲避；又因\Person{科斯蒂利乌斯}公然展示出那\OriginalTerm{敌对者}{adversary}的形象（\BibleRef{得后 2:4}）；为使我们的言语不至于冗长，我们最好满足于圣经的权威，并让我们所有人都服从圣经针对其他异端、尤其是针对这一异端所给的命令。那命令如下：\BibleQuote{离开，离开，从那里出来，不要沾不洁之物；从他们中间出来，要洁净自己，你们抬主器皿的人啊！}（\BibleRef{依 52:11}）这足以教导我们众人，倘若有人曾被他们欺骗，就该从他们中间出来，像从\Place{索多玛}出来一样，且不要再回到他们当中，免得遭受罗特妻子的命运；倘若有人自始便保守纯洁，未受这邪恶异端的玷污，他就可在基督内夸耀说：\Quote{我们未曾向一位外邦的神伸手；我们未曾崇拜自己手造的工，也未侍奉受造物过于侍奉你——那藉着你的圣言、我们的主耶稣基督、你的\OriginalTerm{独生子}{Only-Begotten Son}创造万有的天主。愿光荣和权能，藉着同一圣言，在圣神内，归于你圣父，至于无穷之世。阿们。}\par

% structure: p0036 source=h2 target=subsection reason=relative-heading-rank; base=h2

\subsection{八十一、第二次抗议书}

在\Place{亚历山大}的天主教会之信众，在至可敬的\Person{亚大纳削}主教治理下，由署名之人公开发表本抗议书。\par

我们已经就针对我们自身及\OriginalTerm{主的殿}{Lord's house}发动的夜间袭击提出过抗议；尽管事实上，鉴于全城皆已知晓这些行径，本无需抗议。因为被发现的遇害者遗体公之于众，在主的殿内发现的弓、箭及其他兵器，也大声宣告了这桩恶行。\par

然而，在我们提出抗议之后，最显赫的统帅\Person{西里亚努斯}竟竭力迫使众人附和他，仿佛既无骚乱发生，也无人死亡（这充分证明这些事并非出于最仁慈的\Person{奥古斯都君士坦提乌斯}皇帝的意愿；因为如果他曾奉命行事，便不会如此畏惧此事的后果）；此外，当我们前去见他，请求他既不要对任何人施行暴力，也不要否认已发生之事时，他竟下令用棍棒殴打身为基督徒的我们；这便再次证明针对教会发动的夜间突袭确有其事：——\par

因此，我们现也呈上这份抗议书，我们中一些人即将前往最虔敬的奥古斯都皇帝那里；我们以全能天主之名，并为最虔敬的\Person{奥古斯都君士坦提乌斯}的救恩之故，恳请埃及总督\Person{马克西穆斯}及各位\OriginalTerm{检察官}{Controllers}，将这一切禀报皇帝的虔诚，并报知最显赫的诸位总督的权威。我们也恳请船主们将这些事四处传扬，传至最虔敬的奥古斯都、各地总督及官长耳中，好让人知道针对教会挑起了一场战争，且在\Person{奥古斯都君士坦提乌斯}治下，\Person{西里亚努斯}使\OriginalTerm{贞女}{virgins}及许多其他人成为\OriginalTerm{殉道者}{martyrs}。\par

在二月的伊杜斯前第五天，即\OriginalTerm{麦基尔月}{Mechir}第十四日，黎明时分，我们正在\OriginalTerm{主的殿}{Lord's house}中守夜祈祷（因次日\OriginalTerm{预备日}{Preparation}有\OriginalTerm{圣体圣事}{communion}要举行）；约半夜时分，最显赫的统帅\Person{西里亚努斯}率领众多军团士兵，手持出鞘的剑、标枪及其他兵器，头戴盔，向我们和教会发动攻击；正当我们祈祷、读经之际，他们破门而入。当暴徒们将门撞开后，他一声令下，有人开始射击；有人呐喊，兵器铿鸣，刀剑在灯光下闪烁；\OriginalTerm{贞女}{virgins}们立时被杀害，许多男子被践踏，随着士兵扑来，人叠人倒下，多人中箭身亡。有些士兵还趁机劫掠，剥去贞女的衣裳，贞女们畏惧遭其触碰甚于死亡。主教仍端坐于宝座，勉励众人祈祷。统帅率众攻击，身边有书记\Person{希拉里}随行，此人在后续事件中的作用得以显明。主教被捉，侥幸未被撕碎，他陷入昏迷，状若死人，从他们中间消失，不知去向。他们急于杀他。当他们见已有许多人死去，便命令士兵移走尸首，匿于视线之外。那些留下的至圣贞女被葬入墓穴，在最虔敬的\Person{君士坦提乌斯}治下获得了\OriginalTerm{殉道}{martyrdom}的荣耀。执事们甚至在主的殿内被鞭打，并被囚禁在那里。\par

事态甚至未止于此：这一切过后，凡乐意的便可随意破门而入，搜索并劫掠其中之物。他们甚至进入了并非所有基督徒都可进入的地方。城市卫戍部队的指挥官\Person{戈尔格尼乌斯}知晓此事，因他当时在场。攻打者们所携甲胄、标枪及刀剑被遗留在\OriginalTerm{主的殿}{Lord's house}中，这极为有力地证明了这场敌意攻击的性质。这些兵械至今仍悬于教会中，使他们无法否认；尽管他们几度派来士兵\Person{迪纳米乌斯}及城防指挥官，欲将其取走，我们均未允许，直至此事已为众人所知。\par

若已下令要我们受迫害，我们都已准备殉道。但若非\OriginalTerm{奥古斯都}{Augustus}的命令，我们请求\Place{埃及}总督\Person{Maximus}及所有城市长官请求他，不再允许他们如此攻击我们。我们也希望这请愿能呈递给他，免得他们企图带其他主教进来：因为我们已抵抗至死，渴望拥有最可敬的\Person{Athanasius}，他是天主从一开始就赐给我们的，按我们祖先的传承；他也是最虔诚的\OriginalTerm{奥古斯都}{Augustus} \Person{Constantius}亲自以信件和誓言派遣给我们的。我们相信，当圣上得知所发生的事，他会大不悦，不会违背他的誓言，反会再次下令，让我们的主教\Person{Athanasius}与我们同在。\par

致在最杰出的\Person{Arbæthion}和\Person{Collianus}执政官任期之后将选出的执政官，于\OriginalTerm{梅切尔月}{Mechir}十七日，即二月\OriginalTerm{望日}{Idus}前一日。\par

\SourceRecord{Translated by M. Atkinson and Archibald Robertson.；Nicene and Post-Nicene Fathers, Second Series；Vol. 4.；Edited by Philip Schaff and Henry Wace.；Buffalo, NY: Christian Literature Publishing Co.,；1892.；<http://www.newadvent.org/fathers/28158.htm>.}
